\documentclass[11pt]{article}

\usepackage[a4paper,margin=2.4cm]{geometry}
\usepackage{amsmath,amssymb,amsfonts,amsthm}
\usepackage{bm}
\usepackage{graphicx}
\usepackage{xcolor}
\usepackage{hyperref}
\hypersetup{colorlinks=true,linkcolor=blue!55!black,citecolor=blue!55!black,urlcolor=blue!55!black}
\usepackage{tikz}
\usetikzlibrary{arrows.meta,decorations.pathmorphing,positioning}
\usepackage{mathtools}
\usepackage{empheq}
\usepackage{slashed}

\allowdisplaybreaks[4]

% ---------------- macros ----------------
\newcommand{\eps}{\varepsilon}
\newcommand{\dd}{\mathrm{d}}
\newcommand{\ip}[1]{\int\!\frac{\dd^3 #1}{(2\pi)^3}}
\newcommand{\dt}[1]{\dd^3 #1\,}
\newcommand{\half}{\tfrac12}
\newcommand{\he}{\mathrm{h.c.}}
\newcommand{\dg}{\dagger}
\newcommand{\ket}[1]{\left|#1\right\rangle}
\newcommand{\bra}[1]{\left\langle#1\right|}
\newcommand{\me}[3]{\left\langle #1 \right| #2 \left| #3 \right\rangle}
\newcommand{\Tr}{\mathrm{Tr}}
\newcommand{\vac}{\ket{0}}
\newcommand{\Lam}{\Lambda}
\newcommand{\dthree}{\delta^{(3)}}
\newcommand{\gl}[3]{g^{#1}_{#2}(#3)}
\newcommand{\qk}[3]{q^{#1}_{#2}(#3)}
\newcommand{\aqk}[3]{\bar q^{#1}_{#2}(#3)}
\newcommand{\Phiv}{\Phi}
\newcommand{\Gam}{\Gamma}
\newcommand{\Vggg}{\mathcal{V}}
\newcommand{\OM}{\Omega}
\newcommand{\OMq}{\OM_{q}}
\newcommand{\Hqqg}{H_{q\bar q g}}
\newcommand{\Hggqq}{H_{ggq\bar q}}
\newcommand{\Hinst}{H_{\rm inst}}
\newcommand{\Hg}{H_{g}}
\newcommand{\Hggg}{H_{ggg}}
\newcommand{\Hgggg}{H_{gggg}}
\newcommand{\nord}[1]{\;:\!#1\!:\;}

\title{\bf Soft Quark and Gluon Light-Cone Wave Function and\\ Evolution Operator in the CGC at Next-to-Leading Order}
\author{}
\date{}

\begin{document}
\maketitle
\vspace{-1.4cm}

\begin{abstract}
\noindent
We extend the construction of the soft light-cone wave function (LCWF) and of the associated unitary
evolution operator $\OM$ of a fast moving hadron in the Color Glass Condensate (CGC) framework
to include quark degrees of freedom. Working in light-cone gauge $A^+=0$, within the
Born--Oppenheimer separation between fast valence and soft modes and in the eikonal approximation,
we write the quark sector of $\OM$ as a fully normal ordered series in soft gluon, quark and
antiquark creation and annihilation operators, with coefficients that are operator valued
functionals of the valence colour charge density $\rho^a=\rho^a_g+\rho^a_q$. The coefficients are
fixed by two complementary sets of conditions: the unitarity requirement $\OM^\dg\OM=1$, imposed
order by order in the coupling, which relates the coefficients to one another, and explicit
light-cone perturbation theory (LCPT) calculations of the LCWF for vacuum, one-, two- and
three-gluon and $q\bar q$ incoming states. We obtain all coefficients of $\OM$ through $O(g^2)$.
As an application we diagonalize the soft QCD light-cone Hamiltonian with quarks at orders $g$ and
$g^2$: at $O(g)$ the Hamiltonian is fully diagonalized once the $g\to q\bar q$ and $q\to qg$
structures are included in $\OM$, and at $O(g^2)$ the only surviving correction is the coherent
background field energy $\propto\rho^2$, now built from the full (gluon plus quark) valence colour
charge density, all quark induced corrections being either removable by the $O(g^2)$ part of $\OM$
or proportional to scaleless integrals.
\end{abstract}

\tableofcontents
\bigskip

%=========================================================================
\section{Introduction}
%=========================================================================

In a companion work (referred to below as Paper~I) the soft gluon light-cone wave function and the
corresponding unitary evolution operator $\OM$ of a fast moving hadron were constructed explicitly
through $O(g^2)$ in {\em pure} Yang--Mills theory ($N_f=0$) within the CGC framework. The strategy
there was:
\begin{itemize}
\item[(i)] parameterize $\OM$ as a fully normal ordered series in soft gluon creation and annihilation
operators with operator valued coefficients which are functionals of the valence colour charge
density $\rho$;
\item[(ii)] impose $\OM^\dg\OM=1$ order by order in $g$, which fixes the diagonal and vacuum
coefficients in terms of the others;
\item[(iii)] act with $\OM$ on the vacuum, one-, two- and three-gluon Fock states and match, sector
by sector in Fock space, onto the LCWF components computed directly in light-cone perturbation
theory (LCPT).
\end{itemize}
The purpose of the present paper is to carry out exactly the same programme once quark degrees of
freedom are switched on. Quarks enter in two logically distinct ways:

\begin{enumerate}
\item {\bf The valence sector.} The colour charge density which acts as the static background for
the soft modes now receives a fundamental representation contribution from the fast valence quarks,
$\rho^a=\rho^a_g+\rho^a_q$. The algebra $[\rho^a(k),\rho^b(p)]=if^{abc}\rho^c(k+p)$ is unchanged,
so all pure gluon results of Paper~I go through verbatim with the enlarged $\rho$; what changes is
the operator content of $\rho$ itself.
\item {\bf The soft sector.} The soft Fock space is enlarged by soft quarks and antiquarks. New
operator structures appear in $\OM$: at $O(g)$ the gluon splitting $g\to q\bar q$ and the emission
$q\to qg$, $\bar q\to\bar q g$; at $O(g^2)$ a large set of structures which we enumerate below.
\end{enumerate}

The quark sector is not merely a technical addition. Single inclusive quark production, quark
contributions to the running of the coupling in the NLO evolution kernel, quark--antiquark
correlations in the dilute--dense regime, and the quark contribution to the NLO JIMWLK/BK kernel
all require the explicit coefficients of the quark part of $\OM$. Moreover, as we show, the $O(g)$
diagonalization of the light-cone Hamiltonian in the presence of quarks {\em fails} unless the
$q\to qg$ and $\bar q\to \bar q g$ structures are retained in $\OM$ --- these are the quark
analogues of the $a^\dg a^\dg a$ (three-gluon) structure of Paper~I.

The paper is organized as follows. Section~\ref{sec:setup} fixes conventions, presents the quark
sector of the light-cone QCD Hamiltonian, the Born--Oppenheimer/eikonal separation and the
enlarged colour charge density. Section~\ref{sec:lcpt} collects the LCPT rules and all matrix
elements needed below. Section~\ref{sec:ansatz} gives the complete normal ordered ansatz for the
quark part of $\OM$ and derives the full set of unitarity constraints at $O(g)$ and $O(g^2)$.
Section~\ref{sec:action} works out the action of $\OM$ on the relevant Fock states.
Section~\ref{sec:wf} computes the LCWFs from LCPT. Section~\ref{sec:coeff} collects the resulting
coefficients. Section~\ref{sec:diag} performs the diagonalization at $O(g)$ and $O(g^2)$.
Section~\ref{sec:gen} constructs the Hermitian generator $G$ with $\OM=e^{iG}$.
Section~\ref{sec:concl} concludes. Appendices collect the two-component spinor conventions and
identities, the derivation of the instantaneous vertices, the complete list of matrix elements,
and a summary of the differences with respect to the preliminary notes on which this work is based.

%=========================================================================
\section{Setup and conventions}
\label{sec:setup}
%=========================================================================

\subsection{Light-cone kinematics}

We use the same conventions as in Paper~I. A four vector is written as
$x^\mu=(x^+,x^-,\bm x)$ with $x^\pm=(x^0\pm x^3)/\sqrt2$, and quantization is performed on
surfaces of fixed $x^+$. Momenta are $p^\mu=(p^+,p^-,\bm p)$, and for an on-shell massless
particle
\begin{equation}
p^-=\eps_p\equiv \frac{\bm p^2}{2p^+}\,.
\end{equation}
We use the shorthand
\begin{equation}
\dd^3p\equiv \dd p^+\dd^2 \bm p,\qquad
\ip{p}\equiv\int\!\frac{\dd^3p}{(2\pi)^3},\qquad
\dthree(k-p)\equiv \delta(k^+-p^+)\,\delta^{(2)}(\bm k-\bm p)\,.
\end{equation}
Throughout we work in light-cone gauge $A^+=0$, in which the transverse gluon field $A^i$
($i=1,2$) and the ``good'' component of the quark field are the only dynamical degrees of freedom.
Bold face is dropped whenever no confusion can arise: $k^i$ always denotes a transverse component
and $k^2\equiv\bm k^2$.

For a state consisting of massless quanta with momenta $\{p_a\}$ the light-cone energy is
$E=\sum_a\eps_{p_a}$; we write $\eps_{q\bar q}(l,m)=\eps_l+\eps_m$, $\eps_{gg}(p,q)=\eps_p+\eps_q$
etc.

\subsection{Mode expansions and Fock space}

The transverse gluon field is expanded as
\begin{equation}
A^a_i(x)=\ip{k}\frac{1}{\sqrt{2k^+}}\Big[a^a_i(k)e^{-ik\cdot x}+a^{\dg a}_i(k)e^{ik\cdot x}\Big],
\qquad
\big[a^a_i(k),a^{\dg b}_j(p)\big]=(2\pi)^3\delta^{ab}\delta_{ij}\dthree(k-p).
\label{eq:Aexp}
\end{equation}
For the quark field we use the standard two-component (``$\chi$'') representation of the good
component $\psi_+=\Lambda_+\psi$, $\Lambda_\pm=\tfrac12\gamma^\mp\gamma^\pm$. Writing
$\psi_+=2^{-1/4}\binom{\xi}{0}$ with $\xi$ a two-component spinor in helicity space and
$\alpha,\beta,\dots$ fundamental colour indices, $f,h,\dots$ flavour indices,
\begin{equation}
\xi^{\alpha,f}(x)=\sum_{\lambda=\pm}\ip{k}
\Big[b^{\alpha,f}_\lambda(k)\,\chi_\lambda\,e^{-ik\cdot x}
+d^{\dg\alpha,f}_\lambda(k)\,\chi_{-\lambda}\,e^{ik\cdot x}\Big],
\label{eq:psiexp}
\end{equation}
with the canonical anticommutators
\begin{equation}
\big\{b^{\alpha,f}_\lambda(k),b^{\dg\beta,h}_{\lambda'}(p)\big\}
=\big\{d^{\alpha,f}_\lambda(k),d^{\dg\beta,h}_{\lambda'}(p)\big\}
=(2\pi)^3\delta^{\alpha\beta}\delta^{fh}\delta_{\lambda\lambda'}\dthree(k-p),
\end{equation}
all other (anti)commutators vanishing. Note that, in this convention, the good component of the
quark field carries {\em no} $1/\sqrt{2k^+}$ normalization factor; consequently the longitudinal
momentum factors $1/\sqrt{k^+}$ appearing in the vertices below are always associated with
{\em gluon} legs only. The two-component spinors satisfy
$\chi^\dg_\lambda\chi_{\lambda'}=\delta_{\lambda\lambda'}$ and
$\sum_\lambda\chi_\lambda\chi^\dg_\lambda=\mathbb{1}_{2\times2}$; further identities are collected
in Appendix~\ref{app:spinors}. Flavour indices are carried along but are otherwise inert; sums over
$N_f$ flavours are indicated explicitly where relevant.

The normalized Fock states are
\begin{align}
\ket{\gl{a}{i}{k}}&=\frac{a^{\dg a}_i(k)}{(2\pi)^{3/2}}\vac,\qquad
\ket{\gl{a}{i}{k}\gl{b}{j}{p}}=\frac{a^{\dg a}_i(k)a^{\dg b}_j(p)}{(2\pi)^{3}}\vac,\qquad
\ket{\gl{a}{i}{k}\gl{b}{j}{p}\gl{c}{k}{q}}=\frac{a^{\dg a}_i(k)a^{\dg b}_j(p)a^{\dg c}_k(q)}{(2\pi)^{9/2}}\vac,
\nonumber\\
\ket{\qk{\alpha}{\lambda_1}{l}}&=\frac{b^{\dg\alpha}_{\lambda_1}(l)}{(2\pi)^{3/2}}\vac,\qquad
\ket{\qk{\alpha}{\lambda_1}{l}\,\aqk{\beta}{\lambda_2}{m}}
=\frac{b^{\dg\alpha}_{\lambda_1}(l)d^{\dg\beta}_{\lambda_2}(m)}{(2\pi)^{3}}\vac,
\label{eq:states}
\end{align}
and analogously for mixed states, normalized as
$\langle \gl{b}{j}{p}|\gl{a}{i}{k}\rangle=\delta^{ab}\delta_{ij}\dthree(k-p)$ etc. The soft vacuum
$\vac$ is annihilated by all soft $a$, $b$, $d$ with $\Lam<k^+<\vee$ and is understood to be
tensored with the valence state, $\vac\otimes\ket{v}$.

\subsection{The light-cone QCD Hamiltonian with quarks}

Starting from
$\mathcal L=-\tfrac14 F^a_{\mu\nu}F^{a\,\mu\nu}+\bar\psi\,(i\slashed D)\,\psi$ and eliminating the
constrained components $A^-$ and $\psi_-$ in $A^+=0$ gauge one obtains
$H=H_0+H_{\rm int}$ with
\begin{equation}
H_0=\ip{k}\Big[\eps_k\,a^{\dg a}_i(k)a^a_i(k)
+\eps_k\,b^{\dg\alpha}_\lambda(k)b^\alpha_\lambda(k)
+\eps_k\,d^{\dg\alpha}_\lambda(k)d^\alpha_\lambda(k)\Big],
\end{equation}
(sums over repeated indices, including $\lambda$ and flavour, are implied) and
\begin{equation}
H_{\rm int}=\underbrace{\Hggg+\Hgggg}_{\text{gluon self interactions}}
+\underbrace{\Hqqg}_{\text{quark--gluon vertex}}
+\underbrace{\Hggqq}_{\text{instantaneous quark exchange}}
+\underbrace{H^{(A)}_{\rm inst}}_{\text{instantaneous gluon exchange}} .
\label{eq:Hint}
\end{equation}
The two three- and four-gluon terms are as in Paper~I. The quark--gluon vertex reads, in two
component notation,
\begin{equation}
\Hqqg=-g\!\int\!\dd x^-\dd^2x\;
\xi^\dg\Big\{2\frac{\partial^i}{\partial^+}A^i_a
-\sigma^iA^i_a\frac{\sigma^j\partial^j}{\partial^+}
-\frac{\sigma^j\partial^j}{\partial^+}\sigma^iA^i_a\Big\}t^a\,\xi\,,
\label{eq:Hqqg}
\end{equation}
which contains, in a single expression, the $g\to q\bar q$ splitting, the emissions $q\to qg$ and
$\bar q\to\bar q g$ and their conjugates. The instantaneous quark exchange (the ``seagull''
generated by the non-dynamical component $\psi_-$) is
\begin{equation}
\Hggqq=g^2\!\int\!\dd x^-\dd^2 x\;
\xi^\dg\,\big(\sigma^iA^i_at^a\big)\frac{1}{2i\partial^+}\big(\sigma^jA^j_bt^b\big)\,\xi\,,
\label{eq:Hggqq}
\end{equation}
and the instantaneous gluon exchange takes the compact current--current form
\begin{equation}
H^{(A)}_{\rm inst}=\frac{g^2}{2}\int\!\dd x^-\dd^2x\;
\mathcal J^a\,\frac{1}{(\partial^+)^2}\,\mathcal J^a\,,\qquad
\mathcal J^a=\underbrace{\rho^a}_{\rm valence}
+\underbrace{f^{abc}A^b_i\partial^+A^c_i}_{\mathcal J^a_g}
+\underbrace{\xi^\dg t^a\xi}_{\mathcal J^a_q}\,.
\label{eq:Hinstglue}
\end{equation}
Equation~\eqref{eq:Hinstglue} is a convenient organizing principle. Its various cross terms
reproduce (and extend) the instantaneous interactions used in Paper~I:
\begin{itemize}
\item $\rho\!\cdot\!\mathcal J_g$ gives $H_{gg-{\rm inst}}$ of Paper~I;
\item $\rho\!\cdot\!\mathcal J_q$ gives the instantaneous interaction of the {\em soft quark
current} with the valence charge; it is the term denoted $H_{q\bar q-{\rm inst}}$ in the
preliminary notes, and produces both the $\vac\to\ket{q\bar q}$ and the $\ket{q}\to\ket{q}$
instantaneous transitions;
\item $\mathcal J_q\!\cdot\!\mathcal J_q$ gives the instantaneous gluon exchange inside a soft
$q\bar q$ pair;
\item $\mathcal J_q\!\cdot\!\mathcal J_g$ gives the instantaneous quark--gluon interaction;
\item $\mathcal J_g\!\cdot\!\mathcal J_g$ together with the genuine four-gluon vertex gives
$\Hgggg+H_{gggg-{\rm inst}}$ of Paper~I.
\end{itemize}

\subsection{Born--Oppenheimer separation, eikonal approximation and the colour charge density}
\label{sec:BO}

As in Paper~I, the modes are split into {\em valence} ($k^+\ge\vee$) and {\em soft}
($\Lam<k^+<\vee$) according to their longitudinal momentum,
\begin{equation}
A^a_i=\mathcal A^a_i+\overline{A}^a_i,\qquad
\xi=\xi_{\rm soft}+\overline{\xi},
\end{equation}
where barred fields denote valence modes. Because $k^+_{\rm val}\gg k^+_{\rm soft}$, the valence
sector is static on the time scale of the soft dynamics; all contributions in which
$1/\partial^+$ acts on a valence field are dropped (eikonal approximation). The entire effect of
the valence sector on the soft modes is then encoded in the valence colour charge density.
{\em With quarks, this density has both an adjoint and a fundamental contribution:}
\begin{equation}
\rho^a(-p)=\rho^a_g(-p)+\rho^a_q(-p),
\label{eq:rho}
\end{equation}
\begin{align}
\rho^a_g(-p)&=-if^{abc}\int_\vee^\infty\!\frac{\dd k^+}{2\pi}\!\int\!\frac{\dd^2k}{(2\pi)^2}\;
\bar a^{\dg b}_j(k)\,\bar a^c_j(k+p),\\
\rho^a_q(-p)&=\int_\vee^\infty\!\frac{\dd k^+}{2\pi}\!\int\!\frac{\dd^2k}{(2\pi)^2}\sum_{\lambda,f}
\Big[\bar b^{\dg\alpha,f}_\lambda(k)\,t^a_{\alpha\beta}\,\bar b^{\beta,f}_\lambda(k+p)
-\bar d^{\dg\alpha,f}_\lambda(k+p)\,t^a_{\beta\alpha}\,\bar d^{\beta,f}_\lambda(k)\Big].
\end{align}
The crucial property, which is all that is used below, is unchanged:
\begin{equation}
\big[\rho^a(k),\rho^b(p)\big]=if^{abc}\rho^c(k+p).
\label{eq:rhoalgebra}
\end{equation}
Consequently {\em every} coefficient of $\OM$ is an operator in the valence Hilbert space and
different colour structures do not commute; operator ordering is tracked carefully throughout.
Because soft and valence modes live in different momentum windows, soft operators
$a,a^\dg,b,b^\dg,d,d^\dg$ commute with $\rho$.

After the mode decomposition, the soft--valence coupling of the gluon field is the same operator
as in Paper~I,
\begin{equation}
\Hg=\int_\Lam^\vee\!\frac{\dd k^+}{2\pi}\!\int\!\frac{\dd^2k}{(2\pi)^2}\;
\frac{g\,k^i}{\sqrt2\,|k^+|^{3/2}}\Big[a^{\dg a}_i(k)\rho^a(-k)+a^a_i(k)\rho^a(k)\Big],
\label{eq:Hgop}
\end{equation}
but now with $\rho$ from Eq.~\eqref{eq:rho}. The soft quark current couples to the valence charge
only through $H^{(A)}_{\rm inst}$ (the $\rho\!\cdot\!\mathcal J_q$ cross term), because the
eikonal approximation forbids the direct absorption of a soft quark by the valence sector: a soft
quark cannot be absorbed by a fast quark without transferring a large $k^+$. This is the first
structural difference with respect to the pure gluon case, and it is the reason why the quark
sector of $\OM$ has {\em no} $O(g)$ contribution proportional to $\rho$.

%=========================================================================
\section{Light-cone perturbation theory and matrix elements}
\label{sec:lcpt}
%=========================================================================

\subsection{Perturbation theory with non-commuting matrix elements}

Since the coefficients of $\OM$ are operators in the valence space, ordinary Rayleigh--Schr\"odinger
perturbation theory must be used in the form in which matrix elements are never commuted past one
another. Writing $H=H_0+H_{\rm int}$, $H_0\ket{n^{(0)}}=E^{(0)}_n\ket{n^{(0)}}$, and setting the
soft vacuum energy to zero, the dressed state associated with the incoming free state $\ket{n}$ is
\begin{equation}
\ket{\Psi_n}=\ket{n}
+\sum_{m\neq n}\ket{m}\frac{\me{m}{H_{\rm int}}{n}}{E_n-E_m}
+\sum_{m,i\neq n}\ket{m}\frac{\me{m}{H_{\rm int}}{i}}{E_n-E_m}\frac{\me{i}{H_{\rm int}}{n}}{E_n-E_i}
+\dots
\label{eq:PT}
\end{equation}
with the matrix elements always written {\em from left to right in the order in which the
interactions act}, i.e.\ the {\em latest} interaction stands leftmost. All energies are light-cone
energies $\eps=\bm p^2/2p^+$. Completeness in intermediate states is understood with the measure
$\sum_m\to\sum_{\rm d.o.f.}\int \dt{p_1}\cdots\dt{p_n}\,\frac{1}{S}$, $S$ being the symmetry factor
of identical gluons in the intermediate state.

\subsection{Vertex functions}

It is convenient to introduce once and for all the transverse--spin structures. For a gluon of
momentum $k$ and polarization index $i$ attached to a fermion line whose momenta immediately to
the left and to the right of the vertex are $P_L$ and $P_R$, define
\begin{equation}
\boxed{\;\Gam^i(k;P_L,P_R)\;\equiv\;\frac{2k^i}{k^+}
-\frac{(\bm\sigma\!\cdot\!\bm P_L)\,\sigma^i}{P_L^+}
-\frac{\sigma^i\,(\bm\sigma\!\cdot\!\bm P_R)}{P_R^+}\;}
\label{eq:Gamma}
\end{equation}
and its spinor matrix element
\begin{equation}
\Phiv^i_{\lambda\lambda'}(k;P_L,P_R)\equiv\chi^\dg_\lambda\,\Gam^i(k;P_L,P_R)\,\chi_{\lambda'}\,.
\end{equation}
All three light-cone quark--gluon vertices are governed by Eq.~\eqref{eq:Gamma}, with the
assignment of $P_L,P_R$ dictated by the direction of the fermion number flow:
\begin{center}
\begin{tabular}{lll}
\hline
process & $P_L$ & $P_R$\\
\hline
$g(k)\to q(l)\,\bar q(m)$ & $l$ (outgoing quark) & $m$ (outgoing antiquark)\\
$q(l)\to q(l')\,g(r)$ & $l'$ (outgoing quark) & $l$ (incoming quark)\\
$\bar q(m)\to \bar q(m')\,g(r)$ & $m$ (incoming antiquark) & $m'$ (outgoing antiquark)\\
\hline
\end{tabular}
\end{center}
For the three-gluon vertex we use the tensor of Paper~I,
\begin{align}
\Vggg^{ijl}(k;p,q)=\Big(p^i-q^i+\frac{q^+-p^+}{k^+}k^i\Big)\delta^{jl}
+\Big(k^j+q^j-\frac{k^++q^+}{p^+}p^j\Big)\delta^{il}
+\Big(\frac{k^++p^+}{q^+}q^l-p^l-k^l\Big)\delta^{ij}.
\label{eq:Vggg}
\end{align}

\subsection{Matrix elements}
\label{sec:me}

The complete list used below is (colour indices: $a,b,c,d,e$ adjoint, $\alpha,\beta$ fundamental;
flavour is conserved along the quark line and is suppressed unless needed):

\paragraph{Pure gauge sector (from Paper~I).}
\begin{align}
\me{\gl{a}{i}{k}}{\Hg}{0}&=\frac{g\,k^i\,\rho^a(-k)}{4\pi^{3/2}|k^+|^{3/2}},
\label{eq:M1}\\[2pt]
\me{\gl{c}{l}{q}\gl{b}{j}{p}}{\Hg}{\gl{a}{i}{k}}&=
\delta^{ac}\delta_{il}\dthree(k-q)\frac{g\,p^j\rho^b(-p)}{4\pi^{3/2}|p^+|^{3/2}}
+\big(b,j,p\leftrightarrow c,l,q\big),
\label{eq:M2}\\[2pt]
\me{\gl{c}{l}{q}\gl{b}{j}{p}}{\Hggg}{\gl{a}{i}{k}}&=
\frac{ig f^{abc}\dthree(k-p-q)}{8\pi^{3/2}\sqrt{k^+p^+q^+}}\;\Vggg^{ijl}(k;p,q),
\label{eq:M3}\\[2pt]
\me{\gl{c}{j}{q}\gl{b}{i}{p}}{H^{(A)}_{\rm inst}}{0}&=
\frac{ig^2f^{abc}(q^+-p^+)\delta_{ij}\rho^a(-p-q)}{2(2\pi)^3\sqrt{p^+q^+}\,(p^++q^+)^2},
\label{eq:M4}\\[2pt]
\me{\gl{c}{j}{q}}{H^{(A)}_{\rm inst}}{\gl{b}{i}{p}}&=
\frac{ig^2f^{abc}(p^++q^+)\delta_{ij}\rho^a(p-q)}{2(2\pi)^3\sqrt{p^+q^+}\,(p^+-q^+)^2}.
\label{eq:M5}
\end{align}

\paragraph{Quark sector.} From Eq.~\eqref{eq:Hqqg},
\begin{align}
\me{\qk{\alpha}{\lambda_1}{l}\,\aqk{\beta}{\lambda_2}{m}}{\Hqqg}{\gl{a}{i}{k}}
&=\frac{g\,t^a_{\alpha\beta}}{8\pi^{3/2}\sqrt{k^+}}\,\dthree(k-l-m)\,
\Phiv^i_{\lambda_1\lambda_2}(k;l,m),
\label{eq:Mq1}\\[2pt]
\me{\qk{\alpha'}{\lambda'}{l'}\gl{d}{n}{r}}{\Hqqg}{\qk{\alpha}{\lambda}{l}}
&=\frac{g\,t^d_{\alpha'\alpha}}{8\pi^{3/2}\sqrt{r^+}}\,\dthree(l-l'-r)\,
\Phiv^n_{\lambda'\lambda}(r;l',l),
\label{eq:Mq2}\\[2pt]
\me{\aqk{\beta'}{\lambda'}{m'}\gl{d}{n}{r}}{\Hqqg}{\aqk{\beta}{\lambda}{m}}
&=-\frac{g\,t^d_{\beta\beta'}}{8\pi^{3/2}\sqrt{r^+}}\,\dthree(m-m'-r)\,
\Phiv^n_{\lambda\lambda'}(r;m,m'),
\label{eq:Mq3}
\end{align}
the relative minus sign and the transposed colour matrix in Eq.~\eqref{eq:Mq3} being the usual
consequence of the reversed fermion-number flow along an antiquark line. From
Eq.~\eqref{eq:Hinstglue}, the $\rho\!\cdot\!\mathcal J_q$ cross term,
\begin{align}
\me{\qk{\alpha}{\lambda_1}{l}\,\aqk{\beta}{\lambda_2}{m}}{H^{(A)}_{\rm inst}}{0}
&=\frac{g^2\,t^a_{\alpha\beta}\,\rho^a(-l-m)\,\delta_{\lambda_1\lambda_2}}{8\pi^3\,(l^++m^+)^2},
\label{eq:Mq4}\\[2pt]
\me{\qk{\alpha'}{\lambda'}{l'}}{H^{(A)}_{\rm inst}}{\qk{\alpha}{\lambda}{l}}
&=\frac{g^2\,t^a_{\alpha'\alpha}\,\rho^a(l-l')\,\delta_{\lambda\lambda'}}{8\pi^3\,(l^+-l'^+)^2},
\qquad
\me{\aqk{\beta'}{\lambda'}{m'}}{H^{(A)}_{\rm inst}}{\aqk{\beta}{\lambda}{m}}
=-\frac{g^2\,t^a_{\beta\beta'}\,\rho^a(m-m')\,\delta_{\lambda\lambda'}}{8\pi^3\,(m^+-m'^+)^2},
\label{eq:Mq5}
\end{align}
and the $\mathcal J_q\!\cdot\!\mathcal J_q$ term,
\begin{equation}
\me{\qk{\alpha'}{\lambda_1'}{l'}\aqk{\beta'}{\lambda_2'}{m'}}{H^{(A)}_{\rm inst}}
{\qk{\alpha}{\lambda_1}{l}\aqk{\beta}{\lambda_2}{m}}
=-\frac{g^2\,t^a_{\alpha'\alpha}t^a_{\beta\beta'}\,
\delta_{\lambda_1\lambda_1'}\delta_{\lambda_2\lambda_2'}}
{(2\pi)^3\,(l^+-l'^+)^2}\,\dthree(l+m-l'-m').
\label{eq:Mq6}
\end{equation}
Finally, from Eq.~\eqref{eq:Hggqq}, the instantaneous quark exchange. Denoting by $P^+_{\rm inst}$
the $+$ momentum flowing along the instantaneous fermion line,
\begin{align}
\me{\qk{\alpha}{\lambda_1}{l}\aqk{\beta}{\lambda_2}{m}}{\Hggqq}{\gl{a}{i}{k}\gl{b}{j}{p}}
&=\frac{g^2\,\dthree(k+p-l-m)}{16\pi^3\sqrt{k^+p^+}}\;
\chi^\dg_{\lambda_1}\Big[\frac{(t^at^b)_{\alpha\beta}\,\sigma^i\sigma^j}{l^+-k^+}
+\frac{(t^bt^a)_{\alpha\beta}\,\sigma^j\sigma^i}{l^+-p^+}\Big]\chi_{\lambda_2},
\label{eq:Mq7}\\[2pt]
\me{\qk{\alpha}{\lambda_1}{l}\aqk{\beta}{\lambda_2}{m}\gl{d}{n}{r}}{\Hggqq}{\gl{a}{i}{k}}
&=\frac{g^2\,\dthree(k-l-m-r)}{16\pi^3\sqrt{k^+r^+}}\;
\chi^\dg_{\lambda_1}\Big[\frac{(t^dt^a)_{\alpha\beta}\,\sigma^n\sigma^i}{l^++r^+}
-\frac{(t^at^d)_{\alpha\beta}\,\sigma^i\sigma^n}{m^++r^+}\Big]\chi_{\lambda_2}.
\label{eq:Mq8}
\end{align}
The signs of the two terms in Eq.~\eqref{eq:Mq8} follow from $P^+_{\rm inst}=l^++r^+$ when the
outgoing gluon attaches on the quark side, and $P^+_{\rm inst}=-(m^++r^+)$ when it attaches on the
antiquark side. Details are given in Appendix~\ref{app:inst}.

%=========================================================================
\section{The evolution operator: ansatz, power counting and unitarity}
\label{sec:ansatz}
%=========================================================================

\subsection{Definition of $\OM$}

The evolution operator is the unitary map from free Fock states to the dressed eigenstates of the
full light-cone Hamiltonian,
\begin{equation}
\ket{\Psi_n}=\OM\ket{n},\qquad \OM^\dg\OM=1,\qquad \OM^\dg H\OM=H_{\rm diag}=\sum_n E_n\ket{n}\!\bra{n}.
\label{eq:Omdef}
\end{equation}
It is the adiabatic evolution operator $U(-\infty,0)$ and it simultaneously (i) generates the soft
LCWF and (ii) diagonalizes the soft Hamiltonian order by order in $g$.

Since the quark and gluon operator structures are independent, we split
\begin{equation}
\OM=\OM_{\rm YM}+\OMq\,,
\end{equation}
where $\OM_{\rm YM}$ is the pure gauge operator of Paper~I (with the {\em enlarged} $\rho$ of
Eq.~\eqref{eq:rho} in its coefficients) and $\OMq$ collects all structures containing at least one
soft quark operator, together with the quark induced correction to the purely gluonic
number-preserving structure. The identity is counted once: we write the total normalization as
$N=N_{\rm YM}+N_q-1$, with $N_q\equiv C_0$ the quark contribution.

\subsection{Collective index notation}

To keep the expressions readable we introduce collective labels
\begin{equation}
1\equiv\{p_1,j_1,b_1\},\quad
\ip{1}\equiv\ip{p_1}\sum_{j_1,b_1},\quad a_1\equiv a^{b_1}_{j_1}(p_1),
\end{equation}
for gluons and
\begin{equation}
\kappa\equiv\{l,\lambda,\alpha,f\},\quad
\ip{\kappa}\equiv\ip{l}\sum_{\lambda,\alpha,f},\quad b_\kappa\equiv b^{\alpha,f}_{\lambda}(l);
\qquad
\bar\kappa\equiv\{m,\bar\lambda,\beta,f\},\quad d_{\bar\kappa}\equiv d^{\beta,f}_{\bar\lambda}(m),
\end{equation}
for quarks and antiquarks. Whenever the argument of $\int\!\dd^3\!\cdot/(2\pi)^3$ is a collective
label the discrete sums are understood as above; for an ordinary momentum argument
(e.g.\ $\int\!\dd^3p/(2\pi)^3$) it denotes the plain momentum integral.
A ``$\delta_{12}$'' below denotes the full Kronecker--Dirac symbol
$(2\pi)^3\delta^{b_1b_2}\delta_{j_1j_2}\dthree(p_1-p_2)$, and similarly $\delta_{\kappa\mu}$.

\subsection{The normal ordered ansatz}

The most general fully normal ordered ansatz containing all structures that can contribute through
$O(g^2)$ is
\begin{align}
\OMq=&\;\big(N_q-1\big)
\nonumber\\
&+\ip{1}\ip{\kappa}\ip{\bar\kappa}
\Big[A_3(1;\kappa\bar\kappa)\,b^\dg_\kappa d^\dg_{\bar\kappa}a_1
+A_4(1;\bar\kappa\kappa)\,a^\dg_1 d_{\bar\kappa}b_\kappa\Big]
\nonumber\\
&+\ip{1}\ip{\kappa'}\ip{\kappa}
\Big[F_1(1;\kappa'\kappa)\,a^\dg_1 b^\dg_{\kappa'}b_\kappa
+F_2(1;\kappa\kappa')\,b^\dg_{\kappa}b_{\kappa'}a_1\Big]
\nonumber\\
&+\ip{1}\ip{\bar\kappa'}\ip{\bar\kappa}
\Big[G_1(1;\bar\kappa'\bar\kappa)\,a^\dg_1 d^\dg_{\bar\kappa'}d_{\bar\kappa}
+G_2(1;\bar\kappa\bar\kappa')\,d^\dg_{\bar\kappa}d_{\bar\kappa'}a_1\Big]
\nonumber\\
&+\ip{\kappa}\ip{\bar\kappa}
\Big[A_1(\kappa\bar\kappa)\,b^\dg_\kappa d^\dg_{\bar\kappa}
+A_2(\bar\kappa\kappa)\,d_{\bar\kappa}b_\kappa\Big]
\nonumber\\
&+\ip{1}\ip{2}\,\mathfrak{B}(1,2)\,a^\dg_1a_2
+\ip{\kappa}\ip{\mu}P(\kappa,\mu)\,b^\dg_\kappa b_\mu
+\ip{\bar\kappa}\ip{\bar\mu}\bar P(\bar\kappa,\bar\mu)\,d^\dg_{\bar\kappa}d_{\bar\mu}
\nonumber\\
&+\ip{1}\ip{2}\ip{\kappa}\ip{\bar\kappa}
\Big[B_1(1,2;\kappa\bar\kappa)\,b^\dg_\kappa d^\dg_{\bar\kappa}a_1a_2
+B_2(1,2;\bar\kappa\kappa)\,a^\dg_1a^\dg_2 d_{\bar\kappa}b_\kappa
\nonumber\\
&\hspace{3.2cm}
+B_3(1,2;\kappa\bar\kappa)\,a^\dg_1a_2\,b^\dg_\kappa d^\dg_{\bar\kappa}
+B_4(1,2;\bar\kappa\kappa)\,a^\dg_1a_2\,d_{\bar\kappa}b_\kappa\Big]
\nonumber\\
&+\ip{1}\ip{2}\ip{3}\ip{4}
\Big[C_1(1;2,3,4;\bar\kappa\kappa)\,a^\dg_4a^\dg_3a^\dg_2a_1\,d_{\bar\kappa}b_\kappa
+C_2(1,2;3,4;\bar\kappa\kappa)\,a^\dg_4a^\dg_3a_2a_1\,d_{\bar\kappa}b_\kappa
\nonumber\\
&\hspace{2.4cm}
+C_3(1,2,3;4;\kappa\bar\kappa)\,b^\dg_\kappa d^\dg_{\bar\kappa}\,a^\dg_4a_3a_2a_1
+C_4(1,2;3,4;\kappa\bar\kappa)\,b^\dg_\kappa d^\dg_{\bar\kappa}\,a^\dg_4a^\dg_3a_2a_1\Big]
\nonumber\\
&+\ip{1}\ip{2}
\Big[D_1(1,2;\kappa\bar\kappa,\mu\bar\mu)\,
b^\dg_\kappa d^\dg_{\bar\kappa}b^\dg_\mu d^\dg_{\bar\mu}\,a_1a_2
+D_2(1,2;\bar\mu\mu,\bar\kappa\kappa)\,a^\dg_1a^\dg_2\,d_{\bar\mu}b_\mu d_{\bar\kappa}b_\kappa\Big]
\nonumber\\
&+\ip{\kappa}\ip{\bar\kappa}\ip{\mu}\ip{\bar\mu}
D_3(\kappa\bar\kappa;\bar\mu\mu)\,b^\dg_\kappa d^\dg_{\bar\kappa}d_{\bar\mu}b_\mu\,,
\label{eq:ansatz}
\end{align}
where the integrations over the quark labels in the $B$, $C$, $D_1$, $D_2$ terms are implied.
Structures with more creation or annihilation operators than those displayed, as well as
$O(g^2)$ corrections to structures already present at $O(g)$, contribute only beyond $O(g^2)$ in
the matching sector considered and are omitted.

\subsection{Power counting}

Each elementary quark--gluon vertex carries one power of $g$; each insertion of $\rho$ through
$\Hg$ or $H^{(A)}_{\rm inst}$ carries one or two powers, respectively. Hence
\begin{equation}
\underbrace{A_3,\;A_4,\;F_1,\;F_2,\;G_1,\;G_2}_{O(g)}\;;\qquad
\underbrace{N_q-1,\;A_1,\;A_2,\;\mathfrak B,\;P,\;\bar P,\;B_{1\ldots4},\;C_{1\ldots4},\;
D_{1,2,3}}_{O(g^2)}\,.
\label{eq:counting}
\end{equation}
Two remarks are in order.

\paragraph{(i) No $O(g)$ structure proportional to $\rho$.} In the pure gauge case the
$O(g)$ part of $\OM$ contained the coherent (Weizs\"acker--Williams) coefficient
$A^b_j(p)\propto g\,p^j\rho^b(-p)/(\sqrt{p^+}p^2)$, which is $O(g)$ and $O(\rho)$. In the quark
sector no such structure exists: creating a soft $q\bar q$ pair out of the valence background
requires at least the emission of a gluon by $\rho$ followed by $g\to q\bar q$, i.e.\ $O(g^2)$,
or a single insertion of the instantaneous term \eqref{eq:Mq4}, which is also $O(g^2)$. This is
why $A_1$, and not $A_3$, is the quark analogue of the coherent field, and why it appears only at
NLO.

\paragraph{(ii) The $q\to qg$ structures are mandatory.} The structures $F_{1,2}$ and $G_{1,2}$ do
not appear in the preliminary version of the ansatz. They are, however, $O(g)$ and they are
generated by the very same vertex \eqref{eq:Hqqg} that produces $A_3$. Without them the $O(g)$
diagonalization of Section~\ref{sec:diag} fails: the matrix elements of $\Hqqg$ connecting
$\ket{q\bar q}$ to $\ket{q\bar q g}$ would remain uncancelled. They are the quark analogues of the
$a^\dg a^\dg a$ (three-gluon) coefficient $C$ of Paper~I.

\subsection{Unitarity: general structure}

Expanding
\begin{equation}
\OM=1+\OM_1+\OM_2+O(g^3),\qquad \OM_1=O(g),\quad\OM_2=O(g^2),
\end{equation}
the condition $\OM^\dg\OM=1$ gives
\begin{equation}
\boxed{\;\OM_1^\dg=-\OM_1\;},\qquad
\boxed{\;\OM_2+\OM_2^\dg=-\OM_1^\dg\OM_1=\OM_1^2\;}.
\label{eq:unit}
\end{equation}
These are exact operator identities in the {\em joint} soft $\otimes$ valence Hilbert space; since
the coefficients do not commute, the ordering in $\OM_1^2$ must be kept. It is important to note
that Eq.~\eqref{eq:unit} determines only the {\em Hermitian part} of $\OM_2$; the anti-Hermitian
part is fixed by the matching to the LCWF. This is exactly the split used in Paper~I.

With the $O(g)$ relations already imposed, the anti-Hermitian $\OM_1$ reads
\begin{align}
\OM_1&=\OM_1{}^{\rm YM}+\OM_1{}^{q},\nonumber\\
\OM_1{}^{\rm YM}&=\ip{1}\Big[A_1a^\dg_1-A^\dg_1a_1\Big]
+\ip{1}\ip{2}\ip{3}\Big[C(1;2,3)\,a^\dg_3a^\dg_2a_1-C^\dg(1;2,3)\,a^\dg_3a_2a_1\Big],
\nonumber\\
\OM_1{}^{q}&=\ip{1}\ip{\kappa}\ip{\bar\kappa}\Big[A_3(1;\kappa\bar\kappa)b^\dg_\kappa d^\dg_{\bar\kappa}a_1
-A_3^\dg(1;\kappa\bar\kappa)a^\dg_1 d_{\bar\kappa}b_\kappa\Big]
\nonumber\\
&\quad+\ip{1}\ip{\kappa'}\ip{\kappa}\Big[F_1(1;\kappa'\kappa)a^\dg_1b^\dg_{\kappa'}b_\kappa
-F_1^\dg(1;\kappa'\kappa)\,b^\dg_\kappa b_{\kappa'}a_1\Big]
\nonumber\\
&\quad+\ip{1}\ip{\bar\kappa'}\ip{\bar\kappa}\Big[G_1(1;\bar\kappa'\bar\kappa)a^\dg_1d^\dg_{\bar\kappa'}d_{\bar\kappa}
-G_1^\dg(1;\bar\kappa'\bar\kappa)\,d^\dg_{\bar\kappa}d_{\bar\kappa'}a_1\Big].
\label{eq:Om1}
\end{align}
Here $A_1$ in $\OM_1{}^{\rm YM}$ denotes the Weizs\"acker--Williams coefficient $A^{b_1}_{j_1}(p_1)$
of Paper~I; we keep the notation of Paper~I and trust that no confusion with the quark coefficient
$A_1(\kappa\bar\kappa)$ arises, the arguments always making the distinction clear. To be safe we
write $\mathcal A_1\equiv A^{b_1}_{j_1}(p_1)$ from here on.

\subsection{Unitarity at $O(g)$}

The first of Eqs.~\eqref{eq:unit} immediately gives, structure by structure,
\begin{empheq}[box=\fbox]{align}
A_4(1;\bar\kappa\kappa)&=-\big[A_3(1;\kappa\bar\kappa)\big]^\dg,
\label{eq:U1}\\
F_2(1;\kappa\kappa')&=-\big[F_1(1;\kappa'\kappa)\big]^\dg,
\label{eq:U2}\\
G_2(1;\bar\kappa\bar\kappa')&=-\big[G_1(1;\bar\kappa'\bar\kappa)\big]^\dg,
\label{eq:U3}
\end{empheq}
or, restoring all indices,
\begin{align}
A_4{}^{b;\beta\alpha}_{j;\lambda_2\lambda_1}(p;m,l)&=
-\Big[A_3{}^{b;\alpha\beta}_{j;\lambda_1\lambda_2}(p;l,m)\Big]^\dg,\qquad
F_2{}^{b;\alpha\alpha'}_{j;\lambda\lambda'}(p;l,l')=
-\Big[F_1{}^{b;\alpha'\alpha}_{j;\lambda'\lambda}(p;l',l)\Big]^\dg .
\end{align}
These are the quark counterparts of the pure gauge relations
$A^\dg_{1j}(p)=-A_{2j}(p)$ and $C_2=-C_1^\dg$ of Paper~I. Note that Eqs.~\eqref{eq:U1}--\eqref{eq:U3}
correct the relations quoted in the preliminary notes, which related $A_3$ to $A_1$ and $A_4$ to
$A_2$: those pairs carry {\em different} numbers of gluon operators and therefore cannot be related
by Hermitian conjugation. The correct pairings are $(A_1,A_2)$ and $(A_3,A_4)$, with the first pair
constrained only at $O(g^2)$ (see Eq.~\eqref{eq:U5} below).

\subsection{Unitarity at $O(g^2)$}

Projecting $\OM_2+\OM_2^\dg=\OM_1^2$ onto the independent normal ordered structures of
Eq.~\eqref{eq:ansatz} yields the complete set of relations listed below. In each case the
left-hand side is $[\OM_2]_{\rm structure}+[\OM_2^\dg]_{\rm structure}$ and the right-hand side is
the corresponding projection of $\OM_1^2$, normal ordered with the fermionic contractions
\begin{equation}
d_{\bar\kappa}b_\kappa\,b^\dg_\mu d^\dg_{\bar\mu}
=\delta_{\kappa\mu}\delta_{\bar\kappa\bar\mu}
-\delta_{\kappa\mu}\,d^\dg_{\bar\mu}d_{\bar\kappa}
-\delta_{\bar\kappa\bar\mu}\,b^\dg_\mu b_\kappa
+b^\dg_\mu b_\kappa\,d^\dg_{\bar\mu}d_{\bar\kappa}\,.
\label{eq:fercontr}
\end{equation}

\paragraph{Vacuum (normalization).}
The identity component of $\OM_1^2$ receives no contribution from the quark sector: every quark
structure in $\OM_1$ contains at least one soft fermion creation or annihilation operator, and
the only fully contracted products either contain a leftover $a^\dg a$ (from $A_3^\dg A_3$) or a
leftover $b^\dg b$ (from $F_1^\dg F_1$). Hence
\begin{empheq}[box=\fbox]{align}
N_q=C_0=1+O(g^4).
\label{eq:U4}
\end{empheq}
This confirms the result $C_0=1$ quoted in the preliminary notes and identifies its origin: the
vacuum$\to q\bar q$ amplitude $A_1$ starts at $O(g^2)$, so $|A_1|^2=O(g^4)$.

\paragraph{$b^\dg d^\dg$ ($q\bar q$ creation out of the vacuum).}
The only contraction producing $b^\dg_\kappa d^\dg_{\bar\kappa}$ with no leftover gluon operator is
$\big(A_3\,b^\dg d^\dg a\big)\big(\mathcal A\,a^\dg\big)$, giving
\begin{empheq}[box=\fbox]{align}
A_1(\kappa\bar\kappa)+\big[A_2(\bar\kappa\kappa)\big]^\dg
=\ip{1}A_3(1;\kappa\bar\kappa)\,\mathcal A_1 .
\label{eq:U5}
\end{empheq}
This is the quark-sector analogue of $B_1=\tfrac12(-2B^\dg_2+A^\dg A^\dg+\dots)$ of Paper~I: it
relates the $O(g^2)$ pair creation amplitude to the product of the $O(g)$ splitting amplitude
$A_3$ and the $O(g)$ coherent field $\mathcal A$. Note the ordering: $A_3$ stands to the {\em left}
of $\mathcal A$, since the gluon is emitted by $\rho$ {\em before} it splits.

\paragraph{$a^\dg a$ (quark loop correction to gluon propagation).}
\begin{empheq}[box=\fbox]{align}
\mathfrak B(1,2)+\big[\mathfrak B(2,1)\big]^\dg
=-\ip{\kappa}\ip{\bar\kappa}\,A^\dg_3(1;\kappa\bar\kappa)\,A_3(2;\kappa\bar\kappa)\,.
\label{eq:U6}
\end{empheq}
This is the statement that the probability for a soft gluon to remain a gluon is reduced by the
probability of its splitting into a $q\bar q$ pair. It is the object which, upon insertion into
$\OM^\dg H\OM$, produces the quark loop contribution to the soft gluon self-energy.

\paragraph{$b^\dg b$ and $d^\dg d$ (quark and antiquark number preserving).}
\begin{empheq}[box=\fbox]{align}
P(\kappa,\mu)+\big[P(\mu,\kappa)\big]^\dg
&=-\ip{1}\ip{\kappa'}F^\dg_1(1;\kappa'\kappa)\,F_1(1;\kappa'\mu),
\label{eq:U7}\\
\bar P(\bar\kappa,\bar\mu)+\big[\bar P(\bar\mu,\bar\kappa)\big]^\dg
&=-\ip{1}\ip{\bar\kappa'}G^\dg_1(1;\bar\kappa'\bar\kappa)\,G_1(1;\bar\kappa'\bar\mu).
\label{eq:U8}
\end{empheq}

\paragraph{$b^\dg d^\dg a a$ (pair creation from two gluons).}
With $\mathcal S_{12}$ denoting symmetrization in the two annihilated gluon labels,
\begin{empheq}[box=\fbox]{align}
B_1(1,2;\kappa\bar\kappa)+\big[B_2(1,2;\bar\kappa\kappa)\big]^\dg
=-\mathcal S_{12}\Big\{&A_3(1;\kappa\bar\kappa)\,\mathcal A^\dg_2
+\mathcal A^\dg_2\,A_3(1;\kappa\bar\kappa)
\nonumber\\&
+\ip{3}A_3(3;\kappa\bar\kappa)\,C^\dg(3;1,2)\Big\}.
\label{eq:U9}
\end{empheq}

\paragraph{$a^\dg a\, b^\dg d^\dg$ (pair creation with a spectator gluon).}
\begin{empheq}[box=\fbox]{align}
B_3(1,2;\kappa\bar\kappa)+\big[B_4(1,2;\bar\kappa\kappa)\big]^\dg
=\;&A_3(2;\kappa\bar\kappa)\,\mathcal A_1+\mathcal A_1\,A_3(2;\kappa\bar\kappa)
\nonumber\\
&+\ip{3}\Big[A_3(3;\kappa\bar\kappa)\,C(2;3,1)+A_3(3;\kappa\bar\kappa)\,C(2;1,3)\Big]
\nonumber\\
&+\ip{\mu}F_1(1;\kappa\mu)\,A_3(2;\mu\bar\kappa)
+\ip{\bar\mu}G_1(1;\bar\kappa\bar\mu)\,A_3(2;\kappa\bar\mu).
\label{eq:U10}
\end{empheq}

\paragraph{$b^\dg d^\dg d b$ ($q\bar q\to q\bar q$).}
\begin{empheq}[box=\fbox]{align}
D_3(\kappa\bar\kappa;\bar\mu\mu)+\big[D_3(\mu\bar\mu;\bar\kappa\kappa)\big]^\dg
=-\ip{1}\Big\{&A_3(1;\kappa\bar\kappa)A^\dg_3(1;\mu\bar\mu)
\nonumber\\
&+F^\dg_1(1;\kappa\mu)\,G_1(1;\bar\kappa\bar\mu)
+G^\dg_1(1;\bar\kappa\bar\mu)\,F_1(1;\kappa\mu)\Big\}.
\label{eq:U11}
\end{empheq}

\paragraph{$b^\dg d^\dg b^\dg d^\dg aa$ (two pairs from two gluons).}
\begin{empheq}[box=\fbox]{align}
D_1(1,2;\kappa\bar\kappa,\mu\bar\mu)+\big[D_2(1,2;\bar\mu\mu,\bar\kappa\kappa)\big]^\dg
=\mathcal S_{12}\,A_3(1;\kappa\bar\kappa)\,A_3(2;\mu\bar\mu)\,.
\label{eq:U12}
\end{empheq}
The right-hand side is automatically antisymmetric under $\kappa\leftrightarrow\mu$ and
$\bar\kappa\leftrightarrow\bar\mu$ once contracted with the fermionic operators, as it must be.

\paragraph{$C$-type structures.}
\begin{empheq}[box=\fbox]{align}
C_4(1,2;3,4;\kappa\bar\kappa)+\big[C_2(1,2;3,4;\bar\kappa\kappa)\big]^\dg
&=\mathcal S\Big\{A_3(1;\kappa\bar\kappa)\,C(2;3,4)
+\ip{\mu}F_1(3;\kappa\mu)\,B_?\dots\Big\}\Big|_{O(g^2)}
\nonumber\\
&=\mathcal S_{12}\mathcal S_{34}\;A_3(1;\kappa\bar\kappa)\,C(2;3,4),
\label{eq:U13}\\
C_3(1,2,3;4;\kappa\bar\kappa)+\big[C_1(4;1,2,3;\bar\kappa\kappa)\big]^\dg
&=-\mathcal S_{123}\;A_3(1;\kappa\bar\kappa)\,C^\dg(4;2,3)\,.
\label{eq:U14}
\end{empheq}
In Eq.~\eqref{eq:U13} the terms involving $F_1$ or $G_1$ times a $B$ coefficient are $O(g^3)$ and
have been dropped; $\mathcal S$ denotes symmetrization over the indicated gluon labels.

\paragraph{Summary.}
Equations~\eqref{eq:U1}--\eqref{eq:U14} exhaust the unitarity content of $\OM^\dg\OM=1$ through
$O(g^2)$ in the quark sector. They determine
\begin{equation}
A_4,\;F_2,\;G_2,\;N_q\quad\text{completely},
\end{equation}
and fix the Hermitian parts of
\begin{equation}
A_1\;\text{vs.}\;A_2,\quad \mathfrak B,\quad P,\;\bar P,\quad
B_1\;\text{vs.}\;B_2,\quad B_3\;\text{vs.}\;B_4,\quad
C_1\;\text{vs.}\;C_3,\quad C_2\;\text{vs.}\;C_4,\quad
D_1\;\text{vs.}\;D_2,\quad D_3 .
\end{equation}
The remaining (anti-Hermitian) information is supplied by the explicit LCPT matching of
Sections~\ref{sec:action}--\ref{sec:coeff}.

%=========================================================================
\section{Action of $\OM$ on Fock states and the matching rule}
\label{sec:action}
%=========================================================================

\subsection{A universal counting rule}

Because $\OM$ is normal ordered, its action on a free Fock state is elementary: each annihilation
operator either contracts with one of the incoming quanta or annihilates the vacuum, and the
creation operators build up the outgoing Fock component. The powers of $2\pi$ generated in this way
follow a simple universal rule which we now derive once and for all.

Consider a generic term of $\OM$ with $n$ annihilation and $m$ creation operators and coefficient
$K$,
\begin{equation}
\OM\supset \ip{1}\!\cdots\!\ip{n}\ip{1'}\!\cdots\!\ip{m'}\;
K(1',\dots,m';1,\dots,n)\;
\big[\text{$m$ creation ops}\big]\big[\text{$n$ annihilation ops}\big].
\end{equation}
Acting on a normalized $n$-particle state, each contraction produces a factor
$(2\pi)^{3/2}\delta^{(3)}$, which together with the measure $1/(2\pi)^3$ gives $(2\pi)^{-3/2}$ per
annihilated particle; the $m$ creation operators produce $(2\pi)^{3m/2}$ times the normalized
outgoing state, which together with the $m$ measures gives $(2\pi)^{-3m/2}$. Therefore
\begin{equation}
\boxed{\;
\OM\ket{n\text{ particles}}\supset (2\pi)^{-\frac32(n+m)}\!\int\!\dt{k_{1'}}\!\cdots\dt{k_{m'}}\;
K\;\ket{m\text{ particles}}\;}
\label{eq:rule1}
\end{equation}
and, comparing with the LCWF written as
$\ket{\Psi}=\int\dt{k_{1'}}\cdots\dt{k_{m'}}\,\Psi\,\ket{m\text{ particles}}$,
\begin{equation}
\boxed{\;K=(2\pi)^{\frac32(n+m)}\;\Psi\;}
\label{eq:rule2}
\end{equation}
This is the master matching relation used throughout Section~\ref{sec:coeff}. It reproduces all the
$2\pi$ factors of Paper~I: $\mathcal A$ ($n=0,m=1$) carries $(2\pi)^{3/2}$, $B_2$ ($n=0,m=2$) and
$B_3$ ($n=1,m=1$) carry $(2\pi)^3$, $C$ ($n=1,m=2$) carries $(2\pi)^{9/2}$ and $D_1$ ($n=1,m=3$)
carries $(2\pi)^6$.

\subsection{Vacuum incoming state}
\begin{equation}
\OMq\vac=\big(N_q-1\big)\vac
+\frac{1}{(2\pi)^3}\sum\int\dt{l}\dt{m}\,A_1(\kappa\bar\kappa)\,
\ket{\qk{\alpha}{\lambda_1}{l}\aqk{\beta}{\lambda_2}{m}}\,,
\label{eq:Omvac}
\end{equation}
so that $A_1$ is read off directly from the $\vac\to\ket{q\bar q}$ LCWF.

\subsection{Single gluon incoming state}
\begin{align}
\OMq\ket{\gl{a}{i}{k}}=&\;\big(N_q-1\big)\ket{\gl{a}{i}{k}}
+\frac{1}{(2\pi)^3}\int\!\dt{p}\;\mathfrak B(p,k)\ket{\gl{b}{j}{p}}
\nonumber\\
&+\frac{1}{(2\pi)^{9/2}}\sum\int\dt{l}\dt{m}\;A_3(k;\kappa\bar\kappa)\,
\ket{\qk{\alpha}{\lambda_1}{l}\aqk{\beta}{\lambda_2}{m}}
\nonumber\\
&+\underbrace{\frac{1}{(2\pi)^{3}}\sum\int\dt{l}\dt{m}\;A_1(\kappa\bar\kappa)\,
\ket{\gl{a}{i}{k}\qk{\alpha}{\lambda_1}{l}\aqk{\beta}{\lambda_2}{m}}}_{\text{disconnected}}
\nonumber\\
&+\frac{1}{(2\pi)^{6}}\sum\int\dt{r}\dt{l}\dt{m}\;B_3(r,k;\kappa\bar\kappa)\,
\ket{\gl{d}{n}{r}\qk{\alpha}{\lambda_1}{l}\aqk{\beta}{\lambda_2}{m}}+\dots
\label{eq:Omg}
\end{align}
The $A_3$ term fixes the $O(g)$ splitting coefficient; the connected part of the
$\ket{g\,q\bar q}$ component fixes $B_3$ once the disconnected piece (pair produced from the
valence background with the incoming gluon as spectator) has been subtracted.

\subsection{Two-gluon incoming state}
\begin{align}
\OMq\ket{\gl{a}{i}{k}\gl{b}{j}{p}}=&\;
\frac{1}{(2\pi)^{6}}\sum\int\dt{l}\dt{m}\Big[B_1(k,p;\kappa\bar\kappa)+B_1(p,k;\kappa\bar\kappa)\Big]
\ket{\qk{}{}{l}\aqk{}{}{m}}
\nonumber\\
&+\underbrace{\frac{1}{(2\pi)^{9/2}}\sum\int\dt{l}\dt{m}\Big[A_3(k;\kappa\bar\kappa)\ket{\gl{b}{j}{p}q\bar q}
+A_3(p;\kappa\bar\kappa)\ket{\gl{a}{i}{k}q\bar q}\Big]}_{\text{disconnected}}
\nonumber\\
&+\frac{1}{(2\pi)^{9}}\sum\int\dt{r}\dt{s}\dt{l}\dt{m}
\Big[C_4(k,p;r,s;\kappa\bar\kappa)+(k\!\leftrightarrow\! p)\Big]\ket{gg\,q\bar q}
\nonumber\\
&+\frac{1}{(2\pi)^{9}}\sum\int\dt{l}\dt{m}\dt{l'}\dt{m'}
\Big[D_1(k,p;\kappa\bar\kappa,\mu\bar\mu)+(k\!\leftrightarrow\! p)\Big]\ket{q\bar q\,q\bar q}+\dots
\label{eq:Omgg}
\end{align}

\subsection{Three-gluon incoming state}
The only new coefficient is $C_3$, extracted from the $\ket{g\,q\bar q}$ component:
\begin{equation}
\OMq\ket{\gl{a}{i}{k}\gl{b}{j}{p}\gl{c}{k}{q}}\Big|_{\ket{gq\bar q}}=
\frac{1}{(2\pi)^{9}}\sum\int\dt{r}\dt{l}\dt{m}\;
\Big[\sum_{\rm 6\;perm}C_3(k,p,q;r;\kappa\bar\kappa)\Big]\ket{\gl{d}{n}{r}q\bar q}+\text{disc.}
\label{eq:Omggg}
\end{equation}

\subsection{$q\bar q$ incoming state}
\begin{align}
\OMq\ket{\qk{\alpha}{\lambda_1}{l}\aqk{\beta}{\lambda_2}{m}}
=&\;\frac{1}{(2\pi)^{6}}\sum\int\dt{l'}\dt{m'}\;D_3(\mu\bar\mu;\bar\kappa\kappa)\,
\ket{\qk{}{}{l'}\aqk{}{}{m'}}
\nonumber\\
&+\frac{1}{(2\pi)^3}\!\int\!\dt{l'}P(\mu,\kappa)\ket{\qk{}{}{l'}\aqk{\beta}{\lambda_2}{m}}
+\frac{1}{(2\pi)^3}\!\int\!\dt{m'}\bar P(\bar\mu,\bar\kappa)\ket{\qk{\alpha}{\lambda_1}{l}\aqk{}{}{m'}}
\nonumber\\
&-\frac{1}{(2\pi)^{9/2}}\!\int\!\dt{p}\;A^\dg_3(p;\kappa\bar\kappa)\ket{\gl{b}{j}{p}}
\nonumber\\
&+\frac{1}{(2\pi)^{9/2}}\sum\int\dt{r}\Big[\!\int\!\dt{l'}F_1(r;\mu\kappa)\ket{\gl{}{}{r}\qk{}{}{l'}\aqk{\beta}{\lambda_2}{m}}
\nonumber\\&\hspace{2.6cm}
+\!\int\!\dt{m'}G_1(r;\bar\mu\bar\kappa)\ket{\gl{}{}{r}\qk{\alpha}{\lambda_1}{l}\aqk{}{}{m'}}\Big]
+\frac{A_2(\bar\kappa\kappa)}{(2\pi)^{3}}\vac+\dots
\label{eq:Omqq}
\end{align}
The third line, generated by $A_4=-A^\dg_3$, is the $O(g)$ annihilation of the pair into a single
gluon; the fourth and fifth lines, generated by $F_1$ and $G_1$, are the $O(g)$ gluon emissions off
the quark and the antiquark. The new information in this sector is $D_3$, $P$ and $\bar P$, and (in
the vacuum component) $A_2$.

%=========================================================================
\section{Soft light-cone wave functions from LCPT}
\label{sec:wf}
%=========================================================================

We now compute the LCWF components required by the matching. Throughout, $\eps_p=\bm p^2/2p^+$ and
we abbreviate
\begin{equation}
D^{(i)}\equiv E_{\rm in}-E_{\rm intermediate},\qquad D^{(f)}\equiv E_{\rm in}-E_{\rm final}\,.
\end{equation}
It is useful to record the following algebraic identity, which is used repeatedly. Writing, for a
splitting $k\to l+m$,
\begin{equation}
z=\frac{l^+}{k^+},\qquad 1-z=\frac{m^+}{k^+},\qquad
\bm\kappa=\bm l-z\bm k=\frac{m^+\bm l-l^+\bm m}{k^+},
\end{equation}
one finds, using $\{\sigma^i,\sigma^j\}=2\delta^{ij}$,
\begin{empheq}[box=\fbox]{align}
\Gam^i(k;l,m)=\frac{2k^i}{k^+}-\frac{(\bm\sigma\!\cdot\!\bm l)\sigma^i}{l^+}
-\frac{\sigma^i(\bm\sigma\!\cdot\!\bm m)}{m^+}
=\frac{(2z-1)\kappa^i+i\,\epsilon^{ij}\kappa^j\sigma^3}{k^+z(1-z)}\,,
\label{eq:Gammakappa}
\end{empheq}
i.e.\ the vertex depends only on the {\em relative} transverse momentum, as it must. Furthermore
\begin{equation}
\eps_l+\eps_m-\eps_k=\frac{\bm\kappa^2}{2z(1-z)k^+},\qquad
\sum_{i}\mathrm{Tr}\big[\Gam^{i\dg}\Gam^{i}\big]
=\frac{4\bm\kappa^2\big[z^2+(1-z)^2\big]}{(k^+)^2z^2(1-z)^2}\,,
\label{eq:trace}
\end{equation}
the second relation exhibiting the familiar $g\to q\bar q$ splitting function.

\subsection{Vacuum incoming state: $\vac\to\ket{q\bar q}$}
\label{sec:wfvac}

Two diagrams contribute at $O(g^2)$. In the first the valence charge emits a soft gluon of momentum
$k$ through $\Hg$, which then splits into the pair,
\begin{equation}
\ket{\Psi^{(1)}_{q\bar q\rho}}=\sum_{\lambda_1\lambda_2,f}\int\!\dt{l}\dt{m}
\ket{\qk{\alpha}{\lambda_1}{l}\aqk{\beta}{\lambda_2}{m}}\,
\frac{\me{q\bar q}{\Hqqg}{\gl{a}{i}{k}}}{\eps_{q\bar q}}\,
\frac{\me{\gl{a}{i}{k}}{\Hg}{0}}{\eps_k}\Bigg|_{k=l+m},
\end{equation}
which, with Eqs.~\eqref{eq:M1} and \eqref{eq:Mq1}, gives
\begin{equation}
\ket{\Psi^{(1)}_{q\bar q\rho}}=\sum\int\!\dt{l}\dt{m}\ket{q\bar q}\;
\frac{g^2\,t^a_{\alpha\beta}\,\rho^a(-k)\;k^i\Phiv^i_{\lambda_1\lambda_2}(k;l,m)}
{32\pi^3\,(k^+)^2\,\eps_k\,\eps_{q\bar q}}\,.
\label{eq:Psi1vac}
\end{equation}
The second is the instantaneous gluon exchange between the valence charge and the soft quark
current, Eq.~\eqref{eq:Mq4}, entering at first order,
\begin{equation}
\ket{\Psi^{(2)}_{q\bar q\rho}}=-\sum\int\!\dt{l}\dt{m}\ket{q\bar q}\,
\frac{\me{q\bar q}{H^{(A)}_{\rm inst}}{0}}{\eps_{q\bar q}}
=-\sum\int\!\dt{l}\dt{m}\ket{q\bar q}\,
\frac{g^2\,t^a_{\alpha\beta}\rho^a(-k)\,\delta_{\lambda_1\lambda_2}}{8\pi^3(k^+)^2\eps_{q\bar q}}\,.
\label{eq:Psi2vac}
\end{equation}
Using $k=l+m$ and $\{\sigma^i,\sigma^j\}=2\delta^{ij}$,
\begin{equation}
k^i\Gam^i(k;l,m)=\big(4\eps_k-2\eps_{q\bar q}\big)\mathbb 1
-\frac{k^+}{l^+m^+}\,(\bm\sigma\!\cdot\!\bm l)(\bm\sigma\!\cdot\!\bm m),
\end{equation}
so that the $\delta_{\lambda_1\lambda_2}/[(k^+)^2\eps_{q\bar q}]$ piece of Eq.~\eqref{eq:Psi1vac}
{\em exactly cancels} the instantaneous contribution \eqref{eq:Psi2vac}. This cancellation is a
useful check on the relative normalization of $\Hqqg$ and $H^{(A)}_{\rm inst}$. The sum is
\begin{empheq}[box=\fbox]{equation}
\ket{\Psi_{q\bar q\rho}}=-\sum_{\lambda_1\lambda_2,f}\int\!\dt{l}\dt{m}\;
\ket{\qk{\alpha}{\lambda_1}{l}\aqk{\beta}{\lambda_2}{m}}\,
\frac{g^2\,t^a_{\alpha\beta}\rho^a(-k)}{8\pi^3}
\left[\frac{\delta_{\lambda_1\lambda_2}}{k^+\bm k^2}
+\frac{\chi^\dg_{\lambda_1}(\bm\sigma\!\cdot\!\bm l)(\bm\sigma\!\cdot\!\bm m)\chi_{\lambda_2}}
{2l^+m^+\bm k^2\,\eps_{q\bar q}}\right]
\label{eq:Psivacfinal}
\end{empheq}
with $k=l+m$ and
$(\bm\sigma\!\cdot\!\bm l)(\bm\sigma\!\cdot\!\bm m)=\bm l\!\cdot\!\bm m+i\epsilon^{ij}l^im^j\sigma^3$.
This is the quark analogue of the Weizs\"acker--Williams field: the soft $q\bar q$ cloud of the
valence colour charge. Unlike the gluon case it appears first at $O(g^2)$.

\subsection{Single gluon incoming state}

\subsubsection{$\ket{g}\to\ket{q\bar q}$ at $O(g)$}
This is the elementary splitting,
\begin{equation}
\ket{\Psi^a_i(k)_{q\bar q}}=\sum\int\!\dt{l}\dt{m}\ket{q\bar q}\,
\frac{\me{q\bar q}{\Hqqg}{\gl{a}{i}{k}}}{\eps_k-\eps_l-\eps_m}
=\sum\int\!\dt{l}\dt{m}\ket{q\bar q}\,
\frac{g\,t^a_{\alpha\beta}\dthree(k-l-m)\,\Phiv^i_{\lambda_1\lambda_2}(k;l,m)}
{8\pi^{3/2}\sqrt{k^+}\,\big(\eps_k-\eps_l-\eps_m\big)}\,.
\label{eq:Psigqq}
\end{equation}

\subsubsection{$\ket{g}\to\ket{g\,q\bar q}$ at $O(g^2)$}
\label{sec:gqqg}
Six connected topologies contribute. We label the outgoing quark $(l,\lambda_1,\alpha)$, antiquark
$(m,\lambda_2,\beta)$ and gluon $(r,n,d)$, and write
$D^{(f)}=\eps_k-\eps_l-\eps_m-\eps_r$.

\paragraph{(a) Splitting followed by emission off the quark.} With intermediate quark momentum
$p=l+r$ and $D^{(i)}_a=\eps_k-\eps_p-\eps_m$,
\begin{equation}
\Psi^{(a)}=\frac{g^2\,(t^dt^a)_{\alpha\beta}}{64\pi^3\sqrt{r^+k^+}}\;
\dthree(k-l-m-r)\;
\frac{\chi^\dg_{\lambda_1}\Gam^n(r;l,p)\,\Gam^i(k;p,m)\,\chi_{\lambda_2}}{D^{(f)}D^{(i)}_a}\,,
\end{equation}
where the sum over the intermediate quark helicity has been performed with
$\sum_\lambda\chi_\lambda\chi^\dg_\lambda=\mathbb 1$.

\paragraph{(b) Splitting followed by emission off the antiquark.} With $m'=m+r$ and
$D^{(i)}_b=\eps_k-\eps_l-\eps_{m'}$,
\begin{equation}
\Psi^{(b)}=-\frac{g^2\,(t^at^d)_{\alpha\beta}}{64\pi^3\sqrt{r^+k^+}}\;
\dthree(k-l-m-r)\;
\frac{\chi^\dg_{\lambda_1}\Gam^i(k;l,m')\,\Gam^n(r;m',m)\,\chi_{\lambda_2}}{D^{(f)}D^{(i)}_b}\,.
\end{equation}

\paragraph{(c) Gluon splitting $g\to gg$ followed by $g\to q\bar q$.} With $k'=l+m$ and
$D^{(i)}_c=\eps_k-\eps_{k'}-\eps_r$,
\begin{equation}
\Psi^{(c)}=\frac{i g^2 f^{ab'd}\,t^{b'}_{\alpha\beta}}{64\pi^3\,k'^+\sqrt{k^+r^+}}\;
\dthree(k-l-m-r)\;
\frac{\Vggg^{ij'n}(k;k',r)\,\Phiv^{j'}_{\lambda_1\lambda_2}(k';l,m)}{D^{(f)}D^{(i)}_c}\,.
\end{equation}
This contribution is absent from the preliminary notes; it is the only $\rho$-independent
$O(g^2)$ topology in which the non-Abelian three-gluon vertex feeds the quark channel, and it is
essential for the $O(g^2)$ diagonalization.

\paragraph{(d) Emission off the valence charge, two time orderings.} Here the emitted gluon
$(r,n,d)$ comes from $\rho$ and the incoming gluon splits; the two orderings give
\begin{equation}
\Psi^{(d)}=\frac{g^2\,t^a_{\alpha\beta}\,\rho^d(-r)\;r^n\,\Phiv^i_{\lambda_1\lambda_2}(k;l,m)}
{32\pi^3\sqrt{k^+}(r^+)^{3/2}\;D^{(f)}}
\left[\frac{1}{\eps_k-\eps_l-\eps_m}-\frac{1}{\eps_r}\right]\dthree(k-l-m)\,,
\end{equation}
the first term corresponding to ``split then emit'' and the second to ``emit then split''
($D^{(i)}=-\eps_r$ in that case). Note that $t^a$ is a numerical matrix acting on the soft quark
colour indices and therefore commutes with $\rho^d$; no ordering ambiguity arises at this order.

\paragraph{(e) Instantaneous quark exchange.} From Eq.~\eqref{eq:Mq8} at first order,
\begin{equation}
\Psi^{(e)}=\frac{g^2\,\dthree(k-l-m-r)}{16\pi^3\sqrt{k^+r^+}\,D^{(f)}}\;
\chi^\dg_{\lambda_1}\left[\frac{(t^dt^a)_{\alpha\beta}\sigma^n\sigma^i}{l^++r^+}
-\frac{(t^at^d)_{\alpha\beta}\sigma^i\sigma^n}{m^++r^+}\right]\chi_{\lambda_2}\,,
\end{equation}
also absent from the preliminary notes.

\paragraph{Disconnected piece.} Finally there is the contribution in which the pair is produced
from the valence background while the incoming gluon is a spectator; it reproduces exactly
$\ket{\gl{a}{i}{k}}\otimes\ket{\Psi_{q\bar q\rho}}$ of Eq.~\eqref{eq:Psivacfinal} and is accounted
for by the $A_1$ term of Eq.~\eqref{eq:Omg}. It must be subtracted before extracting $B_3$.

\subsection{Two-gluon incoming states}

\subsubsection{$\ket{gg}\to\ket{q\bar q}$}
Four topologies contribute. Writing $D^{(f)}=\eps_k+\eps_p-\eps_l-\eps_m$:

\paragraph{(i) Fusion through the three-gluon vertex.} The two incoming gluons merge into a gluon
$q'=k+p$ which then splits; with $D^{(i)}_1=\eps_k+\eps_p-\eps_{q'}$,
\begin{equation}
\Psi^{(i)}_{q\bar q}=-\frac{ig^2f^{abc'}\,t^{c'}_{\alpha\beta}}{64\pi^3\,q'^+\sqrt{k^+p^+}}\;
\dthree(k+p-l-m)\;
\frac{\Vggg^{k'ij}(q';k,p)\,\Phiv^{k'}_{\lambda_1\lambda_2}(q';l,m)}{D^{(f)}D^{(i)}_1}\,.
\end{equation}

\paragraph{(ii),(iii) One gluon absorbed by the valence charge.} One of the two incoming gluons,
say $(k,i,a)$, is absorbed by $\rho$ while the other splits. The two time orderings give
\begin{equation}
\Psi^{(ii+iii)}_{q\bar q}=
\frac{g^2\,t^b_{\alpha\beta}\,\rho^a(k)\;k^i\,\Phiv^j_{\lambda_1\lambda_2}(p;l,m)}
{32\pi^3\sqrt{p^+}(k^+)^{3/2}\,D^{(f)}}
\left[\frac{1}{\eps_p-\eps_l-\eps_m}+\frac{1}{\eps_k}\right]\dthree(p-l-m)
\;+\;\big(k,i,a\leftrightarrow p,j,b\big).
\end{equation}

\paragraph{(iv) Instantaneous quark exchange.} From Eq.~\eqref{eq:Mq7},
\begin{equation}
\Psi^{(iv)}_{q\bar q}=\frac{g^2\,\dthree(k+p-l-m)}{16\pi^3\sqrt{k^+p^+}\,D^{(f)}}\;
\chi^\dg_{\lambda_1}\left[\frac{(t^at^b)_{\alpha\beta}\sigma^i\sigma^j}{l^+-k^+}
+\frac{(t^bt^a)_{\alpha\beta}\sigma^j\sigma^i}{l^+-p^+}\right]\chi_{\lambda_2}\,.
\end{equation}

\subsubsection{$\ket{gg}\to\ket{q\bar q\,q\bar q}$}
Each incoming gluon splits independently into a pair. With
$D^{(f)}=\eps_k+\eps_p-\eps_l-\eps_m-\eps_{l'}-\eps_{m'}$ and the two time orderings,
\begin{align}
\Psi_{q\bar q q\bar q}=&\;\frac{g^2\,t^a_{\alpha\beta}t^b_{\alpha'\beta'}}{64\pi^3\sqrt{k^+p^+}}\;
\dthree(k-l-m)\dthree(p-l'-m')\,
\Phiv^i_{\lambda_1\lambda_2}(k;l,m)\,\Phiv^j_{\lambda_1'\lambda_2'}(p;l',m')
\nonumber\\
&\times\frac{1}{D^{(f)}}
\left[\frac{1}{\eps_k+\eps_p-\eps_l-\eps_m-\eps_{p}}+\frac{1}{\eps_k+\eps_p-\eps_k-\eps_{l'}-\eps_{m'}}\right]
\nonumber\\
=&\;\frac{g^2\,t^a_{\alpha\beta}t^b_{\alpha'\beta'}}{64\pi^3\sqrt{k^+p^+}}\;
\dthree(k-l-m)\dthree(p-l'-m')\,
\frac{\Phiv^i(k;l,m)\,\Phiv^j(p;l',m')}{D^{(f)}}
\left[\frac{1}{\eps_k-\eps_l-\eps_m}+\frac{1}{\eps_p-\eps_{l'}-\eps_{m'}}\right],
\label{eq:Psiqqqq}
\end{align}
plus the interchange $(k,i,a)\leftrightarrow(p,j,b)$ combined with
$(\kappa\bar\kappa)\leftrightarrow(\mu\bar\mu)$. The factorized structure of the last line is the
statement that at this order the two pairs are produced independently; genuinely connected
four-quark correlations first arise at $O(g^4)$ (or through the valence averaging over $\rho$).

\subsubsection{$\ket{gg}\to\ket{gg\,q\bar q}$}
Two topologies contribute at $O(g^2)$: one incoming gluon splits into $q\bar q$ while the other
splits into two gluons, and the ordering in which one gluon first splits into two gluons and one of
them subsequently produces the pair. Denoting the outgoing gluons $(s,m,e)$ and $(u,n,f)$ and the
pair $(l,\lambda_1,\alpha)$, $(m,\lambda_2,\beta)$,
\begin{align}
\Psi^{(1)}_{ggq\bar q}&=\frac{ig^2f^{bfe}\,t^a_{\alpha\beta}}{64\pi^3\sqrt{k^+p^+s^+u^+}}
\dthree(k-l-m)\dthree(p-s-u)\;
\frac{\Vggg^{jnm}(p;u,s)\;\Phiv^i_{\lambda_1\lambda_2}(k;l,m)}{D^{(f)}\,\big(\eps_k-\eps_l-\eps_m\big)},
\nonumber\\
\Psi^{(2)}_{ggq\bar q}&=-\frac{ig^2f^{aec'}\,t^{c'}_{\alpha\beta}}{64\pi^3\,q'^+\sqrt{k^+s^+}}\;
\dthree(k-s-q')\dthree(q'-l-m)\,\delta^{bf}\delta_{jn}\dthree(u-p)\,
\frac{\Vggg^{imk'}(k;s,q')\,\Phiv^{k'}(q';l,m)}{D^{(f)}\big(\eps_k-\eps_s-\eps_{q'}\big)},
\end{align}
with $D^{(f)}=\eps_k+\eps_p-\eps_s-\eps_u-\eps_l-\eps_m$, plus $k\leftrightarrow p$.

\subsection{Three-gluon incoming state: $\ket{ggg}\to\ket{g\,q\bar q}$}
Two topologies contribute. With incoming gluons $(k,i,a)$, $(p,j,b)$, $(q,k,c)$, outgoing gluon
$(s,m,e)$ and pair $(l,m)$, and
$D^{(f)}=\eps_k+\eps_p+\eps_q-\eps_s-\eps_l-\eps_{m}$,
\begin{align}
\Psi^{(1)}_{gq\bar q}&=-\frac{ig^2f^{ebc}\,t^a_{\alpha\beta}}{64\pi^3\sqrt{k^+p^+q^+s^+}}\;
\dthree(s-p-q)\dthree(k-l-m)\;
\frac{\Vggg^{mjk}(s;p,q)\,\Phiv^i_{\lambda_1\lambda_2}(k;l,m)}
{D^{(f)}\big(\eps_k-\eps_l-\eps_{m}\big)}+\text{5 perm.},
\nonumber\\
\Psi^{(2)}_{gq\bar q}&=-\frac{ig^2f^{ebc}\,t^{a}_{\alpha\beta}}{64\pi^3\,\sqrt{p^+q^+s^+}\sqrt{k^+}}\;
\dthree(s-p-q)\dthree(k-l-m)\;
\frac{\Vggg^{mjk}(s;p,q)\,\Phiv^{i}(k;l,m)}
{D^{(f)}\big(\eps_k+\eps_p+\eps_q-\eps_k-\eps_{s}\big)}+\text{5 perm.},
\end{align}
where the six permutations of the three incoming gluon labels are understood.

\subsection{$q\bar q$ incoming state}

\subsubsection{$\ket{q\bar q}\to\ket{q\bar q}$}
Three topologies contribute at $O(g^2)$: (A) $s$-channel annihilation into a gluon which
re-splits; (B) $t$-channel one-gluon exchange between the quark and the antiquark, in its two time
orderings; (C) instantaneous gluon exchange, Eq.~\eqref{eq:Mq6}. With incoming $(l,\lambda_1,\alpha)$,
$(m,\lambda_2,\beta)$, outgoing $(l',\lambda_1',\alpha')$, $(m',\lambda_2',\beta')$,
$k=l+m=l'+m'$ and $D^{(f)}=\eps_l+\eps_m-\eps_{l'}-\eps_{m'}$,
\begin{align}
\Psi^{(A)}_{q\bar q}&=\frac{g^2\,t^a_{\alpha'\beta'}t^a_{\beta\alpha}}{64\pi^3 k^+}
\dthree(l+m-l'-m')\;
\frac{\Phiv^i_{\lambda_1'\lambda_2'}(k;l',m')\;
\big[\Phiv^i_{\lambda_1\lambda_2}(k;l,m)\big]^{*}}
{D^{(f)}\;\big(\eps_l+\eps_m-\eps_k\big)},
\label{eq:PsiA}\\
\Psi^{(B)}_{q\bar q}&=-\frac{g^2\,t^d_{\alpha'\alpha}t^d_{\beta\beta'}}{64\pi^3\,r^+}
\dthree(l+m-l'-m')\;
\frac{\Phiv^n_{\lambda_1'\lambda_1}(r;l',l)\,\Phiv^n_{\lambda_2\lambda_2'}(r;m,m')}{D^{(f)}}
\left[\frac{1}{\eps_l-\eps_{l'}-\eps_r}+\frac{1}{\eps_m-\eps_{m'}-\eps_r}\right],
\label{eq:PsiB}\\
\Psi^{(C)}_{q\bar q}&=\frac{g^2\,t^a_{\alpha'\alpha}t^a_{\beta\beta'}\,
\delta_{\lambda_1\lambda_1'}\delta_{\lambda_2\lambda_2'}}
{(2\pi)^3(l^+-l'^+)^2\;D^{(f)}}\dthree(l+m-l'-m'),
\label{eq:PsiC}
\end{align}
with $r=l-l'=m'-m$ the exchanged gluon momentum. Diagram (A) is the one computed in the
preliminary notes; (B) and (C) are new and are required for consistency: without them the
$O(g^2)$ diagonalization of the $q\bar q$ sector does not close.

\subsubsection{$\ket{q}\to\ket{q}$ and $\ket{\bar q}\to\ket{\bar q}$}
At $O(g^2)$ a soft quark interacts with the valence background either instantaneously,
Eq.~\eqref{eq:Mq5}, or by emitting a gluon that is subsequently absorbed by $\rho$ (and the reverse
ordering),
\begin{align}
\Psi_{q\to q}=&\;
\frac{g^2\,t^a_{\alpha'\alpha}\rho^a(l-l')\,\delta_{\lambda\lambda'}}{8\pi^3(l^+-l'^+)^2(\eps_l-\eps_{l'})}
\nonumber\\
&+\frac{g^2\,t^d_{\alpha'\alpha}\,\rho^d(l'-l)\;r^n\,\Phiv^n_{\lambda'\lambda}(r;l',l)}
{32\pi^3\sqrt{r^+}\,|r^+|^{3/2}\,(\eps_l-\eps_{l'})}
\left[\frac{1}{\eps_l-\eps_{l'}-\eps_r}+\frac{1}{\eps_r}\right]_{r=l-l'},
\label{eq:Psiqq}
\end{align}
plus the self-energy topology (emission and reabsorption by the same quark), which is proportional
to the scaleless integral $\int\dd^2\kappa/\bm\kappa^2$ and therefore vanishes in dimensional
regularization. The antiquark expression is obtained by $t^a_{\alpha'\alpha}\to-t^a_{\alpha\alpha'}$
and the corresponding relabelling.

\subsection{Normalization}
\label{sec:norm}
The quark induced correction to the normalization of the soft vacuum vanishes at $O(g^2)$, since
$\ket{\Psi_{q\bar q\rho}}$ is already $O(g^2)$ and its norm is $O(g^4)$; this is
Eq.~\eqref{eq:U4}. The normalization of the {\em one gluon} state, on the other hand, receives an
$O(g^2)$ quark contribution from the $q\bar q$ component \eqref{eq:Psigqq}. Using
Eqs.~\eqref{eq:Gammakappa}--\eqref{eq:trace} and $\mathrm{tr}(t^at^a)|_{\rm no\;sum}=T_F=\tfrac12$,
\begin{align}
\big\langle\Psi^{a'}_{i'}(k')_{q\bar q}\big|\Psi^{a}_{i}(k)_{q\bar q}\big\rangle
&=\delta^{aa'}\delta_{ii'}\dthree(k-k')\;\frac{g^2N_f}{16\pi^3}
\int_{\Lam/k^+}^{1-\Lam/k^+}\!\!\dd z\,\big[z^2+(1-z)^2\big]\int\!\frac{\dd^2\kappa}{\bm\kappa^2}
\nonumber\\
&\xrightarrow[\ \text{dim.\ reg.}\ ]{}\;0\,,
\label{eq:normqq}
\end{align}
because the transverse integral is scaleless. Note that, in contrast with the gluon loop of
Paper~I, the longitudinal integral $\int_0^1\dd z\,[z^2+(1-z)^2]=2/3$ is finite: the quark
splitting function has no $1/z$ or $1/(1-z)$ singularity, so no $\log(\vee/\Lam)$ is generated.
Consequently
\begin{empheq}[box=\fbox]{equation}
N_q=1-\frac{g^2N_f}{48\pi^3}\int\!\frac{\dd^2\kappa}{\bm\kappa^2}\Big|_{\rm one\;gluon\;sector}=1,
\end{empheq}
consistently with the unitarity result \eqref{eq:U4}.

%=========================================================================
\section{Coefficients of the evolution operator through $O(g^2)$}
\label{sec:coeff}
%=========================================================================

We now apply the master rule \eqref{eq:rule2}, $K=(2\pi)^{\frac32(n+m)}\Psi$, to the wave functions
of Section~\ref{sec:wf}. All coefficients below are operators in the valence Hilbert space; the
ordering displayed is the physical one (latest interaction leftmost).

\subsection{Order $g$}

\paragraph{Gluon splitting $g\to q\bar q$ ($n=1$, $m=2$).}
\begin{empheq}[box=\fbox]{equation}
A_3{}^{b;\alpha\beta}_{j;\lambda_1\lambda_2}(p;l,m)
=\frac{(2\pi)^3\,g\,t^b_{\alpha\beta}}{2\sqrt2\;\sqrt{p^+}}\;
\frac{\dthree(p-l-m)\;\Phiv^j_{\lambda_1\lambda_2}(p;l,m)}{\eps_p-\eps_l-\eps_m}\,.
\label{eq:A3}
\end{empheq}
Equivalently, using Eq.~\eqref{eq:Gammakappa} with $z=l^+/p^+$ and $\bm\kappa=\bm l-z\bm p$,
\begin{equation}
A_3{}^{b;\alpha\beta}_{j;\lambda_1\lambda_2}(p;l,m)
=-\frac{(2\pi)^3\,g\,t^b_{\alpha\beta}}{\sqrt2\;\sqrt{p^+}}\;\dthree(p-l-m)\;
\frac{\chi^\dg_{\lambda_1}\big[(2z-1)\kappa^j+i\epsilon^{jn}\kappa^n\sigma^3\big]\chi_{\lambda_2}}
{\bm\kappa^2}\,,
\end{equation}
a compact form which makes the $1/\bm\kappa$ behaviour and the absence of any $1/z$ singularity
manifest.

\paragraph{Gluon emission off a quark and off an antiquark ($n=1$, $m=2$).}
\begin{empheq}[box=\fbox]{align}
F_1{}^{b;\alpha'\alpha}_{j;\lambda'\lambda}(p;l',l)
&=\frac{(2\pi)^3\,g\,t^b_{\alpha'\alpha}}{2\sqrt2\;\sqrt{p^+}}\;
\frac{\dthree(l-l'-p)\;\Phiv^j_{\lambda'\lambda}(p;l',l)}{\eps_l-\eps_{l'}-\eps_p}\,,
\label{eq:F1}\\
G_1{}^{b;\beta'\beta}_{j;\lambda'\lambda}(p;m',m)
&=-\frac{(2\pi)^3\,g\,t^b_{\beta\beta'}}{2\sqrt2\;\sqrt{p^+}}\;
\frac{\dthree(m-m'-p)\;\Phiv^j_{\lambda\lambda'}(p;m,m')}{\eps_m-\eps_{m'}-\eps_p}\,.
\label{eq:G1}
\end{empheq}
The remaining $O(g)$ coefficients follow from unitarity,
$A_4=-A_3^\dg$, $F_2=-F_1^\dg$, $G_2=-G_1^\dg$, Eqs.~\eqref{eq:U1}--\eqref{eq:U3}.

\subsection{Order $g^2$}

\paragraph{Pair creation from the valence background ($n=0$, $m=2$).}
From Eq.~\eqref{eq:Psivacfinal}, with $(2\pi)^3=8\pi^3$,
\begin{empheq}[box=\fbox]{equation}
A_1{}^{\alpha\beta}_{\lambda_1\lambda_2}(l,m)=-g^2\,t^a_{\alpha\beta}\,\rho^a(-k)
\left[\frac{\delta_{\lambda_1\lambda_2}}{k^+\bm k^2}
+\frac{\chi^\dg_{\lambda_1}(\bm\sigma\!\cdot\!\bm l)(\bm\sigma\!\cdot\!\bm m)\chi_{\lambda_2}}
{2l^+m^+\bm k^2\,(\eps_l+\eps_m)}\right],\qquad k=l+m .
\label{eq:A1}
\end{empheq}
The conjugate coefficient follows from Eq.~\eqref{eq:U5},
$A_2(\bar\kappa\kappa)=\big[\int_1 A_3(1;\kappa\bar\kappa)\mathcal A_1\big]^\dg-A_1^\dg(\kappa\bar\kappa)$.

\paragraph{Consistency check of Eq.~\eqref{eq:U5}.}
It is worth verifying the unitarity relation explicitly, since it tests the relative normalization
of $A_1$, $A_3$ and the Weizs\"acker--Williams coefficient
$\mathcal A^b_j(p)=-\sqrt2\,g\,p^j\rho^b(-p)/(\sqrt{p^+}\bm p^2)$ of Paper~I. Using
Eq.~\eqref{eq:A3},
\begin{equation}
\ip{1}A_3(1;\kappa\bar\kappa)\,\mathcal A_1
=-\frac{g^2}{2}\,\frac{t^a_{\alpha\beta}\rho^a(-k)\;k^j\Phiv^j_{\lambda_1\lambda_2}(k;l,m)}
{k^+\bm k^2\,(\eps_k-\eps_l-\eps_m)}\,.
\end{equation}
On the other hand, computing $A_2$ from the $\ket{q\bar q}\to\vac$ wave function
(second order through a single gluon, plus the first order instantaneous term) one finds that the
two instantaneous contributions cancel between $A_1$ and $A_2^\dg$, while the remaining pieces
combine as
\begin{equation}
A_1+A_2^\dg
=g^2\,t^a_{\alpha\beta}\rho^a(-k)\,k^j\Phiv^j_{\lambda_1\lambda_2}(k;l,m)
\left[\frac{1}{2k^+\bm k^2\eps_{q\bar q}}+\frac{\eps_k}{2k^+\bm k^2\eps_{q\bar q}(\eps_{q\bar q}-\eps_k)}\right]
=\frac{g^2\,t^a\rho^a(-k)\,k^j\Phiv^j}{2k^+\bm k^2(\eps_{q\bar q}-\eps_k)},
\end{equation}
which is precisely $\int_1A_3\mathcal A_1$. Unitarity is therefore satisfied identically, a
nontrivial check on the whole construction.

\paragraph{Quark loop in the gluon sector ($n=m=1$).}
The coefficient $\mathfrak B$ is not obtained by matching: for the purely quark-loop topology
momentum conservation forces the outgoing gluon momentum to equal the incoming one, so that the
LCPT energy denominator degenerates and the contribution is entirely a normalization effect. Its
Hermitian part is fixed by unitarity, Eq.~\eqref{eq:U6},
\begin{equation}
\mathfrak B(1,2)+\big[\mathfrak B(2,1)\big]^\dg
=-\ip{\kappa}\ip{\bar\kappa}A_3^\dg(1;\kappa\bar\kappa)A_3(2;\kappa\bar\kappa),
\end{equation}
and its diagonal value follows from Eq.~\eqref{eq:normqq},
\begin{equation}
\mathfrak B(p,p)\big|_{\rm diag}
=-\frac{g^2N_f}{32\pi^3}\,\delta^{ab}\delta_{ij}\!\int_{\Lam/p^+}^{1-\Lam/p^+}\!\!\dd z\,
\big[z^2+(1-z)^2\big]\!\int\!\frac{\dd^2\kappa}{\bm\kappa^2}\;\longrightarrow\;0 .
\end{equation}
The $\rho$-dependent, momentum non-diagonal part of the gluon-number-preserving coefficient is
supplied by the pure gauge $B_3$ of Paper~I.

\paragraph{Pair creation with a spectator gluon ($n=1$, $m=3$; $K=(2\pi)^6$).}
Collecting the six connected topologies of Section~\ref{sec:gqqg},
\begin{empheq}[box=\fbox]{align}
B_3{}^{da;\alpha\beta}_{ni;\lambda_1\lambda_2}(r,k;l,m)=\;&
\frac{(2\pi)^3g^2}{8\sqrt{r^+k^+}}\,\dthree(k-l-m-r)
\bigg\{
(t^dt^a)_{\alpha\beta}\frac{\chi^\dg\Gam^n(r;l,p)\Gam^i(k;p,m)\chi}{D^{(f)}D^{(i)}_a}
\nonumber\\[2pt]
&\hspace{2.8cm}
-(t^at^d)_{\alpha\beta}\frac{\chi^\dg\Gam^i(k;l,m')\Gam^n(r;m',m)\chi}{D^{(f)}D^{(i)}_b}
\bigg\}
\nonumber\\[2pt]
&+\frac{i(2\pi)^3g^2 f^{ab'd}t^{b'}_{\alpha\beta}}{8\,k'^+\sqrt{k^+r^+}}\,\dthree(k-l-m-r)\,
\frac{\Vggg^{ij'n}(k;k',r)\Phiv^{j'}_{\lambda_1\lambda_2}(k';l,m)}{D^{(f)}D^{(i)}_c}
\nonumber\\[2pt]
&+\frac{(2\pi)^3g^2}{2}\,\frac{\dthree(k-l-m-r)}{\sqrt{k^+r^+}\,D^{(f)}}\,
\chi^\dg_{\lambda_1}\!\left[\frac{(t^dt^a)\sigma^n\sigma^i}{l^++r^+}
-\frac{(t^at^d)\sigma^i\sigma^n}{m^++r^+}\right]\!\chi_{\lambda_2}
\nonumber\\[2pt]
&+\frac{(2\pi)^3g^2}{4}\,t^a_{\alpha\beta}\rho^d(-r)\,
\frac{r^n\Phiv^i_{\lambda_1\lambda_2}(k;l,m)\dthree(k-l-m)}{\sqrt{k^+}(r^+)^{3/2}D^{(f)}}
\left[\frac{1}{\eps_k-\eps_l-\eps_m}-\frac{1}{\eps_r}\right],
\label{eq:B3}
\end{empheq}
with $p=l+r$, $m'=m+r$, $k'=l+m$, $D^{(f)}=\eps_k-\eps_l-\eps_m-\eps_r$,
$D^{(i)}_a=\eps_k-\eps_p-\eps_m$, $D^{(i)}_b=\eps_k-\eps_l-\eps_{m'}$,
$D^{(i)}_c=\eps_k-\eps_{k'}-\eps_r$. The coefficient $B_4$ then follows from
Eq.~\eqref{eq:U10}.

\paragraph{Pair creation from two gluons ($n=2$, $m=2$; $K=(2\pi)^6$).}
\begin{empheq}[box=\fbox]{align}
B_1{}^{ab;\alpha\beta}_{ij;\lambda_1\lambda_2}(k,p;l,m)=\;&
-\frac{i(2\pi)^3g^2f^{abc'}t^{c'}_{\alpha\beta}}{8\,q'^+\sqrt{k^+p^+}}\dthree(k+p-l-m)
\frac{\Vggg^{k'ij}(q';k,p)\Phiv^{k'}_{\lambda_1\lambda_2}(q';l,m)}{D^{(f)}D^{(i)}_1}
\nonumber\\
&+\frac{(2\pi)^3g^2}{2}\frac{\dthree(k+p-l-m)}{\sqrt{k^+p^+}D^{(f)}}
\chi^\dg_{\lambda_1}\!\left[\frac{(t^at^b)\sigma^i\sigma^j}{l^+-k^+}
+\frac{(t^bt^a)\sigma^j\sigma^i}{l^+-p^+}\right]\!\chi_{\lambda_2}
\nonumber\\
&+\frac{(2\pi)^3g^2}{4}\,t^b_{\alpha\beta}\rho^a(k)\,
\frac{k^i\Phiv^j_{\lambda_1\lambda_2}(p;l,m)\dthree(p-l-m)}{\sqrt{p^+}(k^+)^{3/2}D^{(f)}}
\left[\frac{1}{\eps_p-\eps_l-\eps_m}+\frac{1}{\eps_k}\right],
\label{eq:B1}
\end{empheq}
with $q'=k+p$, $D^{(f)}=\eps_k+\eps_p-\eps_l-\eps_m$, $D^{(i)}_1=\eps_k+\eps_p-\eps_{q'}$, plus the
interchange $(k,i,a)\leftrightarrow(p,j,b)$ in the last line. $B_2$ follows from
Eq.~\eqref{eq:U9}.

\paragraph{Two pairs from two gluons ($n=2$, $m=4$; $K=(2\pi)^9$).}
\begin{empheq}[box=\fbox]{equation}
\begin{aligned}
D_1(k,p;\kappa\bar\kappa,\mu\bar\mu)=\;&
\frac{(2\pi)^6g^2\,t^a_{\alpha\beta}t^b_{\alpha'\beta'}}{8\sqrt{k^+p^+}}\,
\dthree(k-l-m)\,\dthree(p-l'-m')\,
\Phiv^i_{\lambda_1\lambda_2}(k;l,m)\,\Phiv^j_{\lambda_1'\lambda_2'}(p;l',m')\\[2pt]
&\times\frac{1}{D^{(f)}}\left[\frac{1}{\eps_k-\eps_l-\eps_m}+\frac{1}{\eps_p-\eps_{l'}-\eps_{m'}}\right]
\end{aligned}
\label{eq:D1}
\end{empheq}
with $D^{(f)}=\eps_k+\eps_p-\eps_l-\eps_m-\eps_{l'}-\eps_{m'}$. Note that
Eq.~\eqref{eq:D1} is exactly the symmetrized product $\mathcal S_{12}A_3(1;\kappa\bar\kappa)
A_3(2;\mu\bar\mu)$ divided by the appropriate energy denominators, in agreement with the structure
of the unitarity relation~\eqref{eq:U12}.

\paragraph{$C$-type coefficients.}
From Section~\ref{sec:wf},
\begin{align}
C_4(k,p;s,u;\kappa\bar\kappa)=\;&
\frac{i(2\pi)^6g^2 f^{bfe}t^a_{\alpha\beta}}{8\sqrt{k^+p^+s^+u^+}}
\dthree(k-l-m)\dthree(p-s-u)
\frac{\Vggg^{jnm}(p;u,s)\Phiv^i_{\lambda_1\lambda_2}(k;l,m)}
{D^{(f)}\big(\eps_k-\eps_l-\eps_m\big)}
\nonumber\\
&-\frac{i(2\pi)^6g^2f^{aec'}t^{c'}_{\alpha\beta}}{8\,q'^+\sqrt{k^+s^+}}
\dthree(k-s-q')\dthree(q'-l-m)\delta^{bf}\delta_{jn}\dthree(u-p)
\nonumber\\
&\hspace{1.2cm}\times
\frac{\Vggg^{imk'}(k;s,q')\Phiv^{k'}_{\lambda_1\lambda_2}(q';l,m)}
{D^{(f)}\big(\eps_k-\eps_s-\eps_{q'}\big)}
\;+\;(k\leftrightarrow p),
\label{eq:C4}\\[6pt]
C_3(k,p,q;s;\kappa\bar\kappa)=\;&
-\frac{i(2\pi)^6g^2f^{ebc}t^a_{\alpha\beta}}{8\sqrt{k^+p^+q^+s^+}}
\dthree(s-p-q)\dthree(k-l-m)\,\Vggg^{mjk}(s;p,q)\,\Phiv^i_{\lambda_1\lambda_2}(k;l,m)
\nonumber\\
&\times\frac{1}{D^{(f)}}\left[\frac{1}{\eps_k-\eps_l-\eps_m}+\frac{1}{\eps_p+\eps_q-\eps_s}\right]
\;+\;\text{5 permutations},
\label{eq:C3}
\end{align}
with the appropriate $D^{(f)}$ in each case. $C_1$ and $C_2$ follow from
Eqs.~\eqref{eq:U13}--\eqref{eq:U14}.

\paragraph{$q\bar q\to q\bar q$ ($n=m=2$; $K=(2\pi)^6$).}
Adding the three topologies \eqref{eq:PsiA}--\eqref{eq:PsiC},
\begin{empheq}[box=\fbox]{align}
D_3(\kappa\bar\kappa;\bar\mu\mu)=\;&
\frac{(2\pi)^3g^2\,t^a_{\alpha'\beta'}t^a_{\beta\alpha}}{8\,k^+}
\dthree(l+m-l'-m')
\frac{\Phiv^i_{\lambda_1'\lambda_2'}(k;l',m')\big[\Phiv^i_{\lambda_1\lambda_2}(k;l,m)\big]^*}
{D^{(f)}\,\big(\eps_l+\eps_m-\eps_k\big)}
\nonumber\\
&-\frac{(2\pi)^3g^2\,t^d_{\alpha'\alpha}t^d_{\beta\beta'}}{8\,r^+}\dthree(l+m-l'-m')\,
\frac{\Phiv^n_{\lambda_1'\lambda_1}(r;l',l)\Phiv^n_{\lambda_2\lambda_2'}(r;m,m')}{D^{(f)}}
\nonumber\\
&\hspace{1.5cm}\times\left[\frac{1}{\eps_l-\eps_{l'}-\eps_r}+\frac{1}{\eps_m-\eps_{m'}-\eps_r}\right]
\nonumber\\
&+\frac{g^2\,t^a_{\alpha'\alpha}t^a_{\beta\beta'}
\delta_{\lambda_1\lambda_1'}\delta_{\lambda_2\lambda_2'}}{(l^+-l'^+)^2\,D^{(f)}}
\dthree(l+m-l'-m'),
\label{eq:D3}
\end{empheq}
with $k=l+m$, $r=l-l'$, $D^{(f)}=\eps_l+\eps_m-\eps_{l'}-\eps_{m'}$.

\paragraph{Soft quark in the valence background ($n=m=1$; $K=(2\pi)^3$).}
From Eq.~\eqref{eq:Psiqq},
\begin{empheq}[box=\fbox]{align}
P^{\alpha'\alpha}_{\lambda'\lambda}(l',l)=\;&
\frac{g^2\,t^a_{\alpha'\alpha}\,\rho^a(l-l')\,\delta_{\lambda\lambda'}}
{(l^+-l'^+)^2\,\big(\eps_l-\eps_{l'}\big)}
\nonumber\\
&+\frac{g^2\,t^d_{\alpha'\alpha}\rho^d(l'-l)\;r^n\Phiv^n_{\lambda'\lambda}(r;l',l)}
{4\,(r^+)^{2}\,\big(\eps_l-\eps_{l'}\big)}
\left[\frac{1}{\eps_l-\eps_{l'}-\eps_r}+\frac{1}{\eps_r}\right]_{r=l-l'},
\label{eq:P}
\end{empheq}
and $\bar P$ is obtained by $t^a_{\alpha'\alpha}\to-t^a_{\alpha\alpha'}$ together with the
antiquark relabelling. These coefficients encode the eikonal colour rotation of a soft quark
traversing the valence field.

\subsection{Summary}
All coefficients of the quark sector of $\OM$ through $O(g^2)$ are now determined:
$A_3,F_1,G_1$ at $O(g)$, Eqs.~\eqref{eq:A3}--\eqref{eq:G1}; $A_4,F_2,G_2$ from unitarity;
$A_1,B_1,B_3,C_3,C_4,D_1,D_3,P,\bar P$ from LCPT matching,
Eqs.~\eqref{eq:A1},\eqref{eq:B3},\eqref{eq:B1},\eqref{eq:D1},\eqref{eq:C4},\eqref{eq:C3},
\eqref{eq:D3},\eqref{eq:P}; and $A_2,B_2,B_4,C_1,C_2,D_2,\mathfrak B,N_q$ from the unitarity
relations \eqref{eq:U4}--\eqref{eq:U14}.

%=========================================================================
\section{Diagonalization of the light-cone Hamiltonian at $O(g)$ and $O(g^2)$}
\label{sec:diag}
%=========================================================================

\subsection{General structure}
Writing $H=H_0+H^{(1)}+H^{(2)}+\dots$ with $H^{(n)}=O(g^n)$ and
$\OM=1+\OM_1+\OM_2+\dots$, and using $\OM_1^\dg=-\OM_1$ and
$\OM_2+\OM_2^\dg=\OM_1^2$, one finds
\begin{align}
\OM^\dg H\OM
&=H_0
\;+\;\Big\{H^{(1)}+\big[H_0,\OM_1\big]\Big\}
\;+\;\Big\{H^{(2)}+H^{(1)}\OM_1+\big[H_0,\OM_2\big]\Big\}+O(g^3).
\label{eq:diagexp}
\end{align}
The $O(g)$ bracket vanishes identically provided $\OM_1$ is chosen as in
Section~\ref{sec:coeff}; the $O(g^2)$ bracket is then reduced by $\OM_2$ to the part which
commutes with $H_0$, i.e.\ to the genuinely energy conserving (``diagonal'') terms.

\subsection{Order $g$}
Write the $O(g)$ interaction in normal ordered operator form. By the same counting rule
\eqref{eq:rule2} applied to $H$ rather than to $\OM$, one has, for each vertex,
$V=(2\pi)^{\frac32(n+m)}\times(\text{matrix element})$, so that
\begin{align}
H^{(1)}=&\;\Hg+\Hggg+\Hqqg
\nonumber\\
=&\;\ip{1}\Big[\mathcal V^{\Hg}_1a^\dg_1+\he\Big]
+\ip{1}\ip{2}\ip{3}\Big[\mathcal V^{ggg}(1;2,3)a^\dg_3a^\dg_2a_1+\he\Big]
\nonumber\\
&+\ip{1}\ip{\kappa}\ip{\bar\kappa}\Big[V_A(1;\kappa\bar\kappa)b^\dg_\kappa d^\dg_{\bar\kappa}a_1+\he\Big]
+\ip{1}\ip{\kappa'}\ip{\kappa}\Big[V_F(1;\kappa'\kappa)a^\dg_1b^\dg_{\kappa'}b_\kappa+\he\Big]
\nonumber\\
&+\ip{1}\ip{\bar\kappa'}\ip{\bar\kappa}\Big[V_G(1;\bar\kappa'\bar\kappa)a^\dg_1d^\dg_{\bar\kappa'}d_{\bar\kappa}+\he\Big].
\end{align}
Using
\begin{align}
\big[H_0,\,b^\dg_\kappa d^\dg_{\bar\kappa}a_1\big]&=\big(\eps_l+\eps_m-\eps_{p_1}\big)\,b^\dg_\kappa d^\dg_{\bar\kappa}a_1,
\nonumber\\
\big[H_0,\,a^\dg_1b^\dg_{\kappa'}b_\kappa\big]&=\big(\eps_{p_1}+\eps_{l'}-\eps_{l}\big)\,a^\dg_1b^\dg_{\kappa'}b_\kappa,
\end{align}
and the explicit coefficients \eqref{eq:A3}--\eqref{eq:G1}, which by construction satisfy
\begin{equation}
A_3=\frac{V_A}{\eps_{p_1}-\eps_l-\eps_m},\qquad
F_1=\frac{V_F}{\eps_{l}-\eps_{l'}-\eps_{p_1}},\qquad
G_1=\frac{V_G}{\eps_{m}-\eps_{m'}-\eps_{p_1}},
\end{equation}
one obtains
\begin{equation}
\big[H_0,\OM_1^q\big]=-\Hqqg\,,
\end{equation}
term by term, while $[H_0,\OM^{\rm YM}_1]=-(\Hg+\Hggg)$ as in Paper~I. Therefore
\begin{empheq}[box=\fbox]{equation}
\OM^\dg\big(H_0+\Hg+\Hggg+\Hqqg\big)\OM=H_0+O(g^2).
\end{empheq}
The Hamiltonian is thus {\em completely} diagonalized at $O(g)$. We stress again that this requires
all three quark structures $A_3$, $F_1$, $G_1$ in $\OM$: the coherent operator alone, or even the
coherent operator supplemented by the three-gluon structure of Paper~I, leaves the $q\to qg$ and
$\bar q\to\bar qg$ matrix elements of $\Hqqg$ uncancelled.

\subsection{Order $g^2$}
At $O(g^2)$ the Hamiltonian is
\begin{equation}
H=H_0+\Hg+\Hggg+\Hqqg+\Hgggg+H^{(A)}_{\rm inst}+\Hggqq\,.
\end{equation}
Inserting the coefficients into Eq.~\eqref{eq:diagexp} and choosing $\OM_2$ to cancel every
structure whose coefficient is multiplied by a nonvanishing energy difference, one is left with the
energy conserving remainder. In the pure gauge sector this reproduces the result of Paper~I,
\begin{equation}
\Big[\OM^\dg H\OM\Big]^{\rm YM}_{O(g^2)}
=-\ip{p}\mathcal A^{\dg b}_j(p)\mathcal A^b_j(p)\,\eps_p
+2\ip{p}\ip{q}\ip{r}\ip{s}\big(\eps_r+\eps_s-\eps_p\big)C^\dg(q;s,r)C(p;r,s)\,a^{\dg}_q a_p,
\end{equation}
the second term vanishing identically. The quark sector adds three new structures:
\begin{align}
\Big[\OM^\dg H\OM\Big]^{q}_{O(g^2)}
=\;&-\ip{1}\ip{2}\ip{\kappa}\ip{\bar\kappa}\big(\eps_l+\eps_m-\eps_{p_1}\big)
A^\dg_3(1;\kappa\bar\kappa)A_3(2;\kappa\bar\kappa)\;a^\dg_1a_2
\nonumber\\
&-\ip{1}\ip{\kappa}\ip{\kappa'}\big(\eps_{l'}+\eps_{p_1}-\eps_l\big)
F^\dg_1(1;\kappa'\kappa)F_1(1;\kappa'\mu)\;b^\dg_\kappa b_\mu
\nonumber\\
&-\ip{1}\ip{\bar\kappa}\ip{\bar\kappa'}\big(\eps_{\bar l'}+\eps_{p_1}-\eps_{\bar l}\big)
G^\dg_1(1;\bar\kappa'\bar\kappa)G_1(1;\bar\kappa'\bar\mu)\;d^\dg_{\bar\kappa}d_{\bar\mu}\,.
\label{eq:quarkdiag}
\end{align}
Each of the three terms is diagonal only for $p_1=p_2$, respectively $\kappa=\mu$; all
non-diagonal parts are removed by $\OM_2$, and, as we now show, the diagonal parts vanish.

\paragraph{The quark loop vanishes.} Using Eqs.~\eqref{eq:Gammakappa}--\eqref{eq:trace}, the first
line of Eq.~\eqref{eq:quarkdiag} evaluated at $p_1=p_2=p$ reads
\begin{equation}
-\ip{\kappa}\ip{\bar\kappa}\big(\eps_l+\eps_m-\eps_p\big)\big|A_3(p;\kappa\bar\kappa)\big|^2
\;\propto\;g^2N_f\int_0^1\!\dd z\,\frac{z^2+(1-z)^2}{2z(1-z)p^+}\int\!\dd^2\kappa\;=\;0,
\end{equation}
since the transverse integral is scaleless. This is the light-cone Hamiltonian counterpart of the
familiar statement that the free gluon self-energy vanishes in dimensional regularization: the
extra factor $(\eps_l+\eps_m-\eps_p)\propto\bm\kappa^2$ turns the logarithmically scaleless
integral of Eq.~\eqref{eq:normqq} into a power-like scaleless one. In the presence of a physical
scale (a target field, a quark mass, or an explicit subtraction) this term is precisely the $N_f$
part of the running of $\alpha_s$ in the soft sector.

\paragraph{The quark and antiquark self-energies vanish.} Identically, the second and third lines
of Eq.~\eqref{eq:quarkdiag} at $\kappa=\mu$ are proportional to $\int\dd^2\kappa$ with no
available scale and hence vanish.

\paragraph{Result.} The only surviving diagonal correction is the coherent background field energy,
\begin{equation}
\ip{p}\mathcal A^{\dg b}_j(p)\mathcal A^b_j(p)\,\eps_p
=\ip{p}\frac{g^2\,\rho^b(p)\rho^b(-p)}{(p^+)^2}\,,
\end{equation}
so that, through $O(g^2)$,
\begin{empheq}[box=\fbox]{equation}
\OM^\dg\,H\,\OM=H_0-\ip{p}\frac{g^2\,\rho^b(p)\rho^b(-p)}{(p^+)^2}\,,
\qquad \rho^b=\rho^b_g+\rho^b_q .
\label{eq:diagfinal}
\end{empheq}

\subsection{Discussion}
Equation~\eqref{eq:diagfinal} is formally identical to the pure Yang--Mills result of Paper~I, but
its content is different in two respects.
\begin{itemize}
\item The colour charge density $\rho$ now includes the {\em valence quark} contribution
$\rho_q$. The Weizs\"acker--Williams background field is generated by the total colour charge, so
that a valence quark source (a proton at moderate $x$, a dipole, a jet) is described on the same
footing as a gluonic source. This is the physically relevant modification of the coherent state
picture.
\item The {\em soft} quark degrees of freedom do not contribute to the $O(g^2)$ ground state
energy. Every quark induced term in $\OM^\dg H\OM$ is either removed by the $O(g^2)$ part of $\OM$
(these are precisely the terms whose coefficients are $B_1,B_3,D_3,P,\bar P,\dots$) or is
proportional to a scaleless integral. The first non-trivial soft-quark contribution to the vacuum
energy therefore arises at $O(g^4)$, through $\rho\to g\to q\bar q\to g\to\rho$.
\end{itemize}
Both statements are consistent with the physical expectation: at the order considered, the soft
quark sector is entirely a {\em fluctuation} on top of the classical gluon background, not a
modification of the background itself.

%=========================================================================
\section{The generator of the evolution}
\label{sec:gen}
%=========================================================================

Writing $\OM=e^{iG}$ with $G=G_1+G_2+\dots$ Hermitian, the expansion
$\OM=1+iG_1+\big(iG_2-\tfrac12G_1^2\big)+\dots$ together with $\OM_1=iG_1$ and
$\OM_2=iG_2-\tfrac12 G_1^2$ gives, using $\OM_2+\OM_2^\dg=\OM_1^2$,
\begin{equation}
G_1=-i\,\OM_1,\qquad G_2=-\frac{i}{2}\big(\OM_2-\OM_2^\dg\big),
\end{equation}
both manifestly Hermitian. The quark sector of the generator is therefore
\begin{align}
G_q=&\;-i\ip{1}\ip{\kappa}\ip{\bar\kappa}
\Big[A_3(1;\kappa\bar\kappa)b^\dg_\kappa d^\dg_{\bar\kappa}a_1-A_3^\dg(1;\kappa\bar\kappa)a^\dg_1d_{\bar\kappa}b_\kappa\Big]
\nonumber\\
&-i\ip{1}\ip{\kappa'}\ip{\kappa}\Big[F_1(1;\kappa'\kappa)a^\dg_1b^\dg_{\kappa'}b_\kappa
-F_1^\dg(1;\kappa'\kappa)b^\dg_\kappa b_{\kappa'}a_1\Big]
\nonumber\\
&-i\ip{1}\ip{\bar\kappa'}\ip{\bar\kappa}\Big[G_1(1;\bar\kappa'\bar\kappa)a^\dg_1d^\dg_{\bar\kappa'}d_{\bar\kappa}
-G_1^\dg(1;\bar\kappa'\bar\kappa)d^\dg_{\bar\kappa}d_{\bar\kappa'}a_1\Big]
\nonumber\\
&-\frac{i}{2}\ip{\kappa}\ip{\bar\kappa}\Big\{\big[A_1-A_2^\dg\big]b^\dg_\kappa d^\dg_{\bar\kappa}
+\big[A_2-A_1^\dg\big]d_{\bar\kappa}b_\kappa\Big\}
\nonumber\\
&-\frac{i}{2}\ip{1}\ip{2}\big[\mathfrak B(1,2)-\mathfrak B^\dg(2,1)\big]a^\dg_1a_2
-\frac{i}{2}\ip{\kappa}\ip{\mu}\big[P(\kappa,\mu)-P^\dg(\mu,\kappa)\big]b^\dg_\kappa b_\mu
\nonumber\\
&-\frac{i}{2}\ip{\bar\kappa}\ip{\bar\mu}\big[\bar P(\bar\kappa,\bar\mu)-\bar P^\dg(\bar\mu,\bar\kappa)\big]
d^\dg_{\bar\kappa}d_{\bar\mu}
\nonumber\\
&-\frac{i}{2}\ip{1}\ip{2}\Big\{\big[B_1-B_2^\dg\big]b^\dg_\kappa d^\dg_{\bar\kappa}a_1a_2
+\big[B_2-B_1^\dg\big]a^\dg_1a^\dg_2 d_{\bar\kappa}b_\kappa
\nonumber\\
&\hspace{2.2cm}
+\big[B_3-B_4^\dg\big]a^\dg_1a_2b^\dg_\kappa d^\dg_{\bar\kappa}
+\big[B_4-B_3^\dg\big]a^\dg_1a_2d_{\bar\kappa}b_\kappa\Big\}
\nonumber\\
&-\frac{i}{2}\Big\{\big[C_3-C_1^\dg\big]b^\dg_\kappa d^\dg_{\bar\kappa}a^\dg_4a_3a_2a_1
+\big[C_4-C_2^\dg\big]b^\dg_\kappa d^\dg_{\bar\kappa}a^\dg_4a^\dg_3a_2a_1+\he\Big\}
\nonumber\\
&-\frac{i}{2}\Big\{\big[D_1-D_2^\dg\big]b^\dg_\kappa d^\dg_{\bar\kappa}b^\dg_\mu d^\dg_{\bar\mu}a_1a_2+\he\Big\}
-\frac{i}{2}\big[D_3-D_3^\dg\big]b^\dg_\kappa d^\dg_{\bar\kappa}d_{\bar\mu}b_\mu\,,
\label{eq:G}
\end{align}
where all integrations are as in Eq.~\eqref{eq:ansatz}. The full generator is $G=G_{\rm YM}+G_q$
with $G_{\rm YM}$ given in Paper~I. Exponentiating $G$ promotes the infinitesimal light-cone boost
into a finite one, resumming repeated soft emissions, absorptions and $q\bar q$ fluctuations into a
single unitary operator on the soft Fock space.

%=========================================================================
\section{Concluding remarks}
\label{sec:concl}
%=========================================================================

We have extended the construction of the soft light-cone wave function and of the associated
unitary evolution operator $\OM$ of a fast moving hadron in the CGC to include quark degrees of
freedom, and we have determined all coefficients of the normal ordered expansion of $\OM$ through
$O(g^2)$.

The main results are the following.
\begin{enumerate}
\item {\bf Complete ansatz.} Equation~\eqref{eq:ansatz} lists all normal ordered structures with
soft quark operators that can contribute through $O(g^2)$. Two classes of structures that are
absent from the naive extension of the pure gauge ansatz are essential: the $O(g)$ emission
coefficients $F_1,G_1$ ($q\to qg$, $\bar q\to\bar q g$), and the $O(g^2)$ number preserving
coefficients $\mathfrak B$, $P$, $\bar P$.
\item {\bf Unitarity relations.} The condition $\OM^\dg\OM=1$ at $O(g)$ gives
$A_4=-A_3^\dg$, $F_2=-F_1^\dg$, $G_2=-G_1^\dg$; at $O(g^2)$ it gives the ten relations
\eqref{eq:U4}--\eqref{eq:U14}, which fix $N_q=1$ and relate $A_1\leftrightarrow A_2$,
$B_1\leftrightarrow B_2$, $B_3\leftrightarrow B_4$, $C_1\leftrightarrow C_3$,
$C_2\leftrightarrow C_4$, $D_1\leftrightarrow D_2$, as well as the Hermitian parts of
$\mathfrak B$, $P$, $\bar P$, $D_3$. Relation \eqref{eq:U5} has been verified explicitly against
the independently computed LCPT wave functions.
\item {\bf Explicit coefficients.} Equations~\eqref{eq:A3}--\eqref{eq:P} give all coefficients in
closed form. Two structural simplifications are worth emphasizing: the vertex depends only on the
relative transverse momentum, Eq.~\eqref{eq:Gammakappa}, and the instantaneous gluon exchange
between the valence charge and the soft quark current cancels exactly against a piece of the
non-instantaneous contribution, leaving the compact result \eqref{eq:A1} for the ``quark
Weizs\"acker--Williams'' amplitude.
\item {\bf Diagonalization.} At $O(g)$ the soft Hamiltonian is completely diagonalized. At $O(g^2)$
the only surviving correction is the coherent background field energy $\propto\rho^2$, now built
from the total valence colour charge $\rho_g+\rho_q$; all soft quark induced corrections either
are removed by $\OM_2$ or vanish as scaleless integrals.
\end{enumerate}

The coefficients derived here are the input required for the computation of quark and gluon
production at NLO in the dilute--dense regime, for the $N_f$ part of the NLO evolution kernel in
the wave function approach, and for quark--antiquark correlations in the soft sector. The
natural next steps are the extension to $O(g^3)$ (needed for single inclusive production at NLO
with quarks in the final state) and the computation of the soft quark number operator expectation
value in the evolved state.

\appendix

%=========================================================================
\section{Two-component spinors: conventions and identities}
\label{app:spinors}
%=========================================================================

We use the chiral representation with $\Lambda_\pm=\tfrac12\gamma^\mp\gamma^\pm$ and
$\psi_\pm=\Lambda_\pm\psi$. The good component is written as
$\psi_+=2^{-1/4}\binom{\xi}{0}$, with $\xi$ a two-component spinor acted on by the Pauli matrices
$\sigma^1,\sigma^2,\sigma^3$. Only $\sigma^{1,2}$ appear contracted with transverse momenta; we use
$\bm\sigma\!\cdot\!\bm p\equiv\sigma^ip^i$ with $i=1,2$. The helicity basis satisfies
\begin{equation}
\sigma^3\chi_\lambda=\lambda\,\chi_\lambda\ (\lambda=\pm1),\qquad
\chi^\dg_\lambda\chi_{\lambda'}=\delta_{\lambda\lambda'},\qquad
\sum_{\lambda}\chi_\lambda\chi^\dg_\lambda=\mathbb 1 .
\end{equation}
The only algebraic identities used in the main text are
\begin{align}
\sigma^i\sigma^j&=\delta^{ij}\mathbb 1+i\epsilon^{ij}\sigma^3\qquad (i,j=1,2;\ \epsilon^{12}=+1),
\label{eq:sigalg}\\
(\bm\sigma\!\cdot\!\bm a)(\bm\sigma\!\cdot\!\bm b)&=\bm a\!\cdot\!\bm b\,\mathbb 1
+i\,\epsilon^{ij}a^ib^j\,\sigma^3,
\qquad
\{\sigma^i,\sigma^j\}=2\delta^{ij}\mathbb 1 .
\end{align}

\paragraph{Derivation of Eq.~\eqref{eq:Gammakappa}.} For a splitting $k\to l+m$ with
$l=zk+\kappa$, $m=(1-z)k-\kappa$, $l^+=zk^+$, $m^+=(1-z)k^+$,
\begin{align}
\Gam^i(k;l,m)&=\frac{1}{k^+}\Big[2k^i-\big(\bm\sigma\!\cdot\!\bm k\big)\sigma^i
-\sigma^i\big(\bm\sigma\!\cdot\!\bm k\big)
-\frac{(\bm\sigma\!\cdot\!\bm\kappa)\sigma^i}{z}+\frac{\sigma^i(\bm\sigma\!\cdot\!\bm\kappa)}{1-z}\Big]
\nonumber\\
&=\frac{1}{k^+}\Big[-\frac{(\bm\sigma\!\cdot\!\bm\kappa)\sigma^i}{z}
+\frac{\sigma^i(\bm\sigma\!\cdot\!\bm\kappa)}{1-z}\Big]
=\frac{(2z-1)\kappa^i+i\epsilon^{ij}\kappa^j\sigma^3}{k^+z(1-z)}\,,
\end{align}
where the first three terms cancel by $\{\sigma^j,\sigma^i\}=2\delta^{ij}$ and the last step uses
Eq.~\eqref{eq:sigalg}. Note that the exact cancellation of the $\bm k$ dependence is a
consequence of gauge invariance (transversality) and provides a strong check on the vertex
\eqref{eq:Gamma}.

\paragraph{Traces.} With $\tau^i\equiv\epsilon^{ij}\kappa^j$ ($\bm\tau\!\cdot\!\bm\kappa=0$,
$\bm\tau^2=\bm\kappa^2$),
\begin{equation}
\mathrm{Tr}\big[\Gam^{i'\dg}\Gam^{i}\big]
=\frac{2\big[\delta^{ii'}\bm\kappa^2-4z(1-z)\kappa^i\kappa^{i'}\big]}{(k^+)^2z^2(1-z)^2},
\end{equation}
which upon angular averaging, $\kappa^i\kappa^{i'}\to\tfrac12\delta^{ii'}\bm\kappa^2$, becomes
\begin{equation}
\mathrm{Tr}\big[\Gam^{i'\dg}\Gam^{i}\big]\Big|_{\rm av}
=\frac{2\,\delta^{ii'}\bm\kappa^2\big[z^2+(1-z)^2\big]}{(k^+)^2z^2(1-z)^2}\,,
\end{equation}
the familiar $g\to q\bar q$ splitting function $z^2+(1-z)^2$ appearing automatically.

%=========================================================================
\section{Instantaneous vertices}
\label{app:inst}
%=========================================================================

In light-cone gauge, two classes of instantaneous interactions arise from eliminating the
non-dynamical fields.

\paragraph{Instantaneous gluon exchange.} Eliminating $A^-$ through the constraint
$\partial^+A^-_a=-\frac{1}{\partial^+}\big(D^{ab}_i\partial^+A^b_i-g\mathcal J^{a}\big)$ generates
the current--current term \eqref{eq:Hinstglue}. Since $\mathcal J^a$ contains the valence charge
$\rho^a$, the soft gluon current $\mathcal J_g^a$ and the soft quark current
$\mathcal J^a_q=\xi^\dg t^a\xi$, all cross terms are obtained at once. In momentum space the
$1/(\partial^+)^2$ becomes $-1/(P^+)^2$ with $P^+$ the longitudinal momentum transferred through
the instantaneous line, which is the origin of the $1/(l^++m^+)^2$ and $1/(l^+-l'^+)^2$ factors in
Eqs.~\eqref{eq:Mq4}--\eqref{eq:Mq6}.

\paragraph{Instantaneous quark exchange.} For massless quarks the constraint equation reads
\begin{equation}
\psi_-=\frac{1}{2i\partial^+}\,\bm\alpha\!\cdot\!\big(i\bm\partial-g\bm A\big)\,\gamma^+\psi_+\,,
\end{equation}
whose $O(g^2)$ piece, substituted back into the Hamiltonian, produces the seagull
\eqref{eq:Hggqq}. Diagrammatically it is a quark line whose propagator has been replaced by the
instantaneous piece $\gamma^+/(2P^+_{\rm inst})$, with $P^+_{\rm inst}$ the $+$ momentum flowing
along the line. The two orderings of the gluon attachments give the two terms in
Eqs.~\eqref{eq:Mq7}--\eqref{eq:Mq8}; the relative minus sign in Eq.~\eqref{eq:Mq8} arises because
in the second ordering the instantaneous line carries $P^+_{\rm inst}=-(m^++r^+)$.

\paragraph{Consistency check.} As shown in Section~\ref{sec:wfvac}, the instantaneous gluon
exchange between $\rho$ and the soft quark current cancels exactly the
$\delta_{\lambda_1\lambda_2}/[(k^+)^2\eps_{q\bar q}]$ part of the non-instantaneous contribution.
Such cancellations between instantaneous and non-instantaneous pieces are a general feature of
light-cone Hamiltonian perturbation theory and are the practical way of checking the relative
normalizations of $\Hqqg$, $\Hggqq$ and $H^{(A)}_{\rm inst}$.

%=========================================================================
\section{Relation to the preliminary notes}
\label{app:diff}
%=========================================================================

For convenience we list the points at which the present treatment differs from, corrects, or
extends the preliminary set of notes on which it is based.

\begin{enumerate}
\item {\bf Unitarity relations at $O(g)$.} The notes quote
$A_3(l,m,p)=-A_1^\dg(p,m,l)$ and $A_4(l,m,p)=-A_2^\dg(p,m,l)$. These pairs carry different numbers
of soft gluon operators ($A_1,A_2$ have none, $A_3,A_4$ have one) and cannot be related by
Hermitian conjugation. The correct statements are $A_4=-A_3^\dg$ at $O(g)$,
Eq.~\eqref{eq:U1}, and the $O(g^2)$ relation $A_1+A_2^\dg=\int_1 A_3\mathcal A_1$,
Eq.~\eqref{eq:U5}, which has been verified explicitly.

\item {\bf Missing $O(g)$ structures.} The gluon emission structures $a^\dg b^\dg b$ and
$a^\dg d^\dg d$ (coefficients $F_1,G_1$) and their conjugates are $O(g)$ and are generated by the
same vertex as $A_3$. They are indispensable for the $O(g)$ diagonalization
(Section~\ref{sec:diag}).

\item {\bf Missing $O(g^2)$ structures.} The number preserving coefficients $\mathfrak B$
($a^\dg a$, quark loop), $P$ ($b^\dg b$) and $\bar P$ ($d^\dg d$) are required for the unitarity
relations to close and for the normalization of the one gluon and one quark states.

\item {\bf The $D_2$ structure.} In the notes the $D_2$ term is written with one creation and one
annihilation gluon operator. To be the conjugate partner of $D_1$ ($b^\dg d^\dg b^\dg d^\dg aa$) it
must carry two gluon {\em creation} operators, $a^\dg a^\dg\,d b\,d b$; this is the form adopted
here.

\item {\bf Energy denominators.} In several places the notes write the denominator of the
$g\to q\bar q$ splitting as $\eps_l+\eps_m$ instead of $\eps_k-\eps_l-\eps_m$ (e.g.\ in the
single-gluon incoming wave function). We use the LCPT rule $E_{\rm in}-E_{\rm intermediate}$
uniformly.

\item {\bf The $q\to qg$ vertex.} The notes write the emission vertex off the quark as
$2r^l/r^+-(\bm\sigma\!\cdot\!\bm p')\sigma^l/p'^+-\sigma^l(\bm\sigma\!\cdot\!\bm q)/q^+$, i.e.\ with
the {\em antiquark} momentum $q$ in the last term. The correct structure involves the momentum of
the incoming quark, $\Gam^l(r;p',p)$, cf.\ the table below Eq.~\eqref{eq:Gamma}.

\item {\bf Missing diagrams.} The following contributions do not appear in the notes and are
included here: (i) $g\to gg$ followed by $g\to q\bar q$ in the $\ket{g}\to\ket{gq\bar q}$ wave
function, topology (c) of Section~\ref{sec:gqqg}; (ii) the instantaneous quark exchange
contributions, Eqs.~\eqref{eq:Mq7}--\eqref{eq:Mq8}, in the $\ket{g}\to\ket{gq\bar q}$ and
$\ket{gg}\to\ket{q\bar q}$ sectors; (iii) the $t$-channel one-gluon exchange and the instantaneous
gluon exchange in $\ket{q\bar q}\to\ket{q\bar q}$, Eqs.~\eqref{eq:PsiB}--\eqref{eq:PsiC}.

\item {\bf Factors of $2\pi$.} The notes carry these somewhat loosely. Here they are fixed once and
for all by the counting rule \eqref{eq:rule2}, which reproduces all the $2\pi$ factors of Paper~I.

\item {\bf The valence charge density.} With quarks, $\rho^a=\rho^a_g+\rho^a_q$, and it is this
enlarged density which enters all coefficients and the final diagonalized Hamiltonian
\eqref{eq:diagfinal}.
\end{enumerate}

\end{document}
